\documentclass[11pt,oneside]{book}

\usepackage[T1]{fontenc}
\usepackage[utf8]{inputenc}
\usepackage{lmodern}
\usepackage{microtype}
\usepackage[a4paper,margin=30mm,headheight=15pt]{geometry}
\usepackage{amsmath,amssymb,amsthm,mathtools}
\usepackage{mathrsfs}
\usepackage{bm}
\usepackage{booktabs,longtable,array,tabularx}
\usepackage{enumitem}
\usepackage{xcolor}
\usepackage{graphicx}
\usepackage{fancyhdr}
\usepackage{algorithm2e}
\usepackage[most]{tcolorbox}
\usepackage{listings}
\usepackage{url}
\usepackage{hyperref}
\usepackage[nameinlink,noabbrev,capitalise]{cleveref}
% Canonical mathematical notation for active FabiusFunction LaTeX documents.
%
% This file is intentionally a plain TeX input rather than a LaTeX package.
% Every standalone document loads it by an explicit path relative to that
% document's directory; included fragments inherit the definitions from their
% root document.  Keep this file independent of document layout and metadata.
%
% Contract:
%   * one command has one mathematical meaning throughout the active corpus;
%   * names encode semantic roles rather than a document-local choice of letter;
%   * visually similar but differently normalized objects have different print;
%   * every command below is defined normatively in the notation catalogue.
%
% The active corpus excludes docs/archive/.  Historical sources may retain
% their original notation only inside an explicitly labelled quotation.

\makeatletter
\@ifundefined{mathbb}{%
  \PackageError{fabius-notation}{The amssymb package must be loaded first}%
    {Load amsmath and amssymb before inputting fabius-notation.tex.}%
}{}
\@ifundefined{operatorname}{%
  \PackageError{fabius-notation}{The amsmath package must be loaded first}%
    {Load amsmath before inputting fabius-notation.tex.}%
}{}
\makeatother

% Version stamp used by the catalogue and by static audits.
\newcommand{\FabiusNotationVersion}{2026-08-31}

% Number systems.  NaturalNumbers is explicitly the nonnegative convention.
\newcommand{\NaturalNumbers}{\mathbb{N}_{0}}
\newcommand{\PositiveIntegers}{\mathbb{N}_{>0}}
\newcommand{\IntegerNumbers}{\mathbb{Z}}
\newcommand{\RationalNumbers}{\mathbb{Q}}
\newcommand{\RealNumbers}{\mathbb{R}}
\newcommand{\ComplexNumbers}{\mathbb{C}}
\newcommand{\ExtendedRealNumbers}{\overline{\mathbb{R}}}

% The three principal Fabius--Rvachev functions.  Subscripts are deliberate:
% a bare F and a calligraphic F have several unrelated uses in the corpus.
\newcommand{\FabiusBounded}{F_{\mathrm{bd}}}
\newcommand{\FabiusGlobal}{F_{\mathrm{gl}}}
\newcommand{\FabiusQuantile}{Q_{F}}
\newcommand{\FabiusClampedQuantile}{Q_{F,\mathrm{cl}}}
\newcommand{\RvachevUp}{\operatorname{up}}

% Constants, elementary operators, and delimiters.
\newcommand{\EulerE}{\mathrm{e}}
\newcommand{\ImaginaryUnit}{\mathrm{i}}
\newcommand{\Differential}{\mathop{}\!\mathrm{d}}
\newcommand{\AbsoluteValue}[1]{\left\lvert #1\right\rvert}
\newcommand{\Norm}[1]{\left\lVert #1\right\rVert}
\newcommand{\Floor}[1]{\left\lfloor #1\right\rfloor}
\newcommand{\Ceiling}[1]{\left\lceil #1\right\rceil}
\newcommand{\FractionalPart}[1]{\left\{#1\right\}}
\newcommand{\PositivePart}[1]{\left(#1\right)_{+}}
\newcommand{\IndicatorOf}[1]{\mathbf{1}_{#1}}
\newcommand{\SetOf}[1]{\left\{#1\right\}}
\newcommand{\InnerProduct}[1]{\left\langle #1\right\rangle}
\newcommand{\CardinalityOf}[1]{\left\lvert #1\right\rvert}
\newcommand{\DefinitionEquals}{\mathrel{\coloneqq}}
\newcommand{\Sign}{\operatorname{sgn}}
\newcommand{\IdentityOperator}{\operatorname{id}}
\newcommand{\SpanOperator}{\operatorname{span}}
\newcommand{\TraceOperator}{\operatorname{tr}}
\newcommand{\RankOperator}{\operatorname{rank}}
\newcommand{\DiagonalOperator}{\operatorname{diag}}
\newcommand{\ResidueOperator}{\operatorname*{Res}}
\newcommand{\LeastCommonMultiple}{\operatorname{lcm}}
\newcommand{\OrderOperator}{\operatorname{ord}}
\newcommand{\PrincipalLogarithm}{\operatorname{Log}}
\newcommand{\FallingFactorial}[2]{\left(#1\right)^{\underline{#2}}}
\newcommand{\RisingFactorial}[2]{\left(#1\right)^{\overline{#2}}}

% Support, topology, convolution, and metric notation.
\newcommand{\SupportOperator}{\operatorname{supp}}
\newcommand{\SupportOf}[1]{\SupportOperator\!\left(#1\right)}
\newcommand{\TopologicalSupportOf}[1]{\operatorname{tsupp}\!\left(#1\right)}
\newcommand{\Convolution}{\mathbin{*}}
\newcommand{\MetricDistance}{\operatorname{dist}}
\newcommand{\TotalVariationDistance}{d_{\mathrm{TV}}}

% Probability.  Equality and convergence in law are not metric distance.
\newcommand{\Probability}{\mathbb{P}}
\newcommand{\Expectation}{\mathbb{E}}
\newcommand{\Variance}{\operatorname{Var}}
\newcommand{\Covariance}{\operatorname{Cov}}
\newcommand{\LawOperator}{\operatorname{Law}}
\newcommand{\LawOf}[1]{\LawOperator\!\left(#1\right)}
\newcommand{\EqualInLaw}{\mathrel{\overset{\mathrm{law}}{=}}}
\newcommand{\ConvergesInLaw}{\mathrel{\xrightarrow{\mathrm{law}}}}

% Fourier and integral transforms.  The printed labels expose normalization.
% FourierTwoPi uses exp(-2*pi*i*x*xi); FourierAngular uses exp(-i*x*omega).
\newcommand{\FourierTwoPi}[1]{\widehat{#1}_{\,2\pi}}
\newcommand{\FourierAngular}[1]{\widehat{#1}_{\,\mathrm{ang}}}
\newcommand{\LaplaceTransformSymbol}{\mathcal{L}}
\newcommand{\MellinTransformSymbol}{\mathcal{M}}
\newcommand{\LaplaceTransformOf}[1]{\LaplaceTransformSymbol\!\left[#1\right]}
\newcommand{\MellinTransformOf}[1]{\MellinTransformSymbol\!\left[#1\right]}
\newcommand{\MomentGeneratingFunction}[1]{M_{#1}}
\newcommand{\CharacteristicFunction}[1]{\varphi_{#1}}

% The two standard sinc functions must never share one symbol.
\newcommand{\SincRad}{\operatorname{sinc}_{\mathrm{rad}}}
\newcommand{\SincPi}{\operatorname{sinc}_{\pi}}

% Binary arithmetic and Thue--Morse notation.
\newcommand{\BinaryDigitSum}{\operatorname{wt}_{2}}
\newcommand{\ThueMorseSignSymbol}{\varepsilon}
\newcommand{\ThueMorseSign}[1]{\varepsilon_{#1}}
\newcommand{\PadicValuation}[1]{\nu_{#1}}
\newcommand{\TwoAdicValuation}{\nu_{2}}

% q-series notation.  Every base is an explicit argument.
\newcommand{\QPochhammer}[3]{\left(#1;#2\right)_{#3}}
\newcommand{\GaussianBinomial}[3]{%
  \genfrac{[}{]}{0pt}{}{#1}{#2}_{#3}%
}
\newcommand{\QInteger}[2]{\left[#1\right]_{#2}}

% Named special functions.
\newcommand{\Polylogarithm}{\operatorname{Li}}
\newcommand{\LambertW}{W}
\newcommand{\GammaFunction}{\Gamma}
\newcommand{\RiemannZeta}{\zeta}

% Asymptotics.  BigO and LittleO are operator tokens for existing formula
% grammar; prefer the At forms so that the limiting regime is printed.
\newcommand{\BigO}{\mathcal{O}}
\newcommand{\LittleO}{o}
\newcommand{\BigOAt}[2]{\mathcal{O}_{#2}\!\left(#1\right)}
\newcommand{\LittleOAt}[2]{o_{#2}\!\left(#1\right)}


\lstset{
  basicstyle=\ttfamily\small,
  columns=fullflexible,
  keepspaces=true,
  breaklines=true,
  breakatwhitespace=false,
  showstringspaces=false,
  frame=single,
  framerule=0.3pt,
  xleftmargin=1em,
  xrightmargin=1em
}

\hypersetup{
  colorlinks=true,
  linkcolor=blue!45!black,
  citecolor=green!35!black,
  urlcolor=blue!55!black,
  pdftitle={Polynomial--Logarithmic Transseries: Algebra, Composition, and Asymptotic Reversion},
  pdfsubject={Polynomial-logarithmic transseries, Lambert W, composition, division, series reversion},
  pdfkeywords={transseries, logarithmic asymptotics, Lambert W, series reversion, composition, Lagrange inversion}
}

\pagestyle{fancy}
\fancyhf{}
\fancyhead[L]{\small\itshape Polynomial--Logarithmic Transseries}
\fancyhead[R]{\small\thepage}
\fancyfoot{}

\setcounter{tocdepth}{2}
\setcounter{secnumdepth}{3}
\allowdisplaybreaks
\setlist[itemize]{topsep=4pt,itemsep=2pt}
\setlist[enumerate]{topsep=4pt,itemsep=2pt}

\definecolor{shadeblue}{RGB}{238,246,252}
\definecolor{shadegray}{RGB}{245,245,245}
\definecolor{shadegreen}{RGB}{239,249,241}
\definecolor{shadered}{RGB}{253,241,241}

\newtheorem{theorem}{Theorem}[chapter]
\newtheorem{proposition}[theorem]{Proposition}
\newtheorem{lemma}[theorem]{Lemma}
\newtheorem{corollary}[theorem]{Corollary}
\theoremstyle{definition}
\newtheorem{definition}[theorem]{Definition}
\newtheorem{example}[theorem]{Example}
\newtheorem{algorithmicprinciple}[theorem]{Algorithmic principle}
\newtheorem{exercise}[theorem]{Exercise}
\theoremstyle{remark}
\newtheorem{remark}[theorem]{Remark}
\newtheorem{warning}[theorem]{Warning}

\newenvironment{abstract}
  {\cleardoublepage\thispagestyle{plain}\begin{center}
   {\Large\bfseries Abstract}\end{center}\begin{quote}\small}
  {\end{quote}\cleardoublepage}

\newtcolorbox{summarybox}{
  breakable,colback=shadeblue,colframe=blue!35!black,
  boxrule=0.5pt,arc=1.5mm,left=2mm,right=2mm,top=1mm,bottom=1mm,
  before skip=8pt,after skip=8pt,fontupper=\small,
  title=Summary,fonttitle=\bfseries\small}
\newtcolorbox{computebox}{
  breakable,colback=shadegreen,colframe=green!35!black,
  boxrule=0.5pt,arc=1.5mm,left=2mm,right=2mm,top=1mm,bottom=1mm,
  before skip=8pt,after skip=8pt,fontupper=\small,
  title=Computation,fonttitle=\bfseries\small}
\newtcolorbox{pitfallbox}{
  breakable,colback=shadered,colframe=red!40!black,
  boxrule=0.5pt,arc=1.5mm,left=2mm,right=2mm,top=1mm,bottom=1mm,
  before skip=8pt,after skip=8pt,fontupper=\small,
  title=Pitfall,fonttitle=\bfseries\small}

\DeclareMathOperator{\lm}{lm}
\DeclareMathOperator{\lc}{lc}
\DeclareMathOperator{\degL}{deg_{\!L}}

\newcommand{\K}{\mathbb K}
\newcommand{\PLpoly}{\PolynomialLogLaurentRing{\K}}
\newcommand{\PLfield}{\IteratedLaurentLogField{\K}}
\newcommand{\D}{\mathcal D}
\newcommand{\EulerDerivation}{\mathcal D_{\mathrm{Eul}}}
\newcommand{\PolynomialLogAsymptotic}{\mathrel{\sim_{\mathrm{PL}}}}
\newcommand{\precdom}{\prec}
\newcommand{\succdom}{\succ}
\newcommand{\compinv}{^{\langle-1\rangle_{\circ}}}
\newcommand{\compose}{\mathbin{\circ}}
\newcommand{\ApproximationOf}[1]{#1_{\mathrm{approx}}}
\newcommand{\trunc}[2]{\left[#1\right]_{\le #2}}
\newcommand{\stirlingone}[2]{\genfrac{[}{]}{0pt}{}{#1}{#2}}
\newcommand{\InverseLogScale}[1]{\varepsilon_{\mathrm{log},#1}}
\newcommand{\LowerBranchParameter}{\varepsilon_{\mathrm{br}}}
\newcommand{\NewtonResidualScale}{\varepsilon_{\mathrm{Newt}}}
\newcommand{\LogarithmOnBranch}[2]{%
  \PrincipalLogarithm_{#1}\!\left(#2\right)%
}

\title{\Huge\bfseries Polynomial--Logarithmic Transseries\\[0.4em]
\Large Algebra, Division, Composition, and Asymptotic Series Reversal\\[1em]
\large with Lambert's \texorpdfstring{$\LambertW$}{W} function as the guiding example}
\author{A self-contained computational monograph}
\date{August 2026}

\begin{document}
\frontmatter
\maketitle

\begin{abstract}
Polynomial--logarithmic transseries are expansions whose successive power
blocks carry polynomial, Laurent-series, or nested-logarithmic coefficients.
The model example is the large-argument expansion of Lambert's function,
\[
 \LambertW(x)=L_1-L_2+\frac{L_2}{L_1}
 +\frac{L_2^2/2-L_2}{L_1^2}+\cdots,
 \qquad L_1=\log x,\quad L_2=\log\log x.
\]
Although such formulas look like ordinary power series with decorative
logarithms, their correct algebra depends on an ordered hierarchy of scales:
$X^{-1}$ is smaller than every fixed negative power of $\log X$, while powers
of $\log X$ remain ordered inside each $X^{-n}$ block.

This article develops a practical and rigorous calculus for these objects.
It defines the polynomial--logarithmic Laurent algebra and its iterated
Laurent-field completion; proves closure and coefficient formulas for
multiplication, reciprocal, and division; derives differentiation formulas
that preserve the class; and gives two complementary constructions of
composition.  The first composes by substituting the generators
$X^{-1}$ and $\log X$; the second is a formally justified Taylor expansion
for near-identity substitutions.

The central part treats series reversal, meaning compositional inversion.
Fixed-point iteration, triangular coefficient recursion, Newton iteration,
and a transserial Lagrange--B\"urmann formula are developed and compared.
For $F(X)=X+A(X)$ the inverse is shown formally to be
\[
 F\compinv(Y)=Y+\sum_{n\ge1}\frac{(-1)^n}{n!}
 \frac{\Differential^{\,n-1}}{\Differential Y^{\,n-1}}A(Y)^n,
\]
whenever the polynomial--logarithmic filtration makes the family summable.
Taking $A(X)=\log X$ yields an all-orders derivation of the Lambert-$\LambertW$
polynomials, their Stirling-number formula, and their integral recurrence.
The discussion also covers complex branches, the lower real branch,
scale-changing inversions such as $Y=cX^\rho(\log X)^\sigma$, nested
logarithms, analytic error transfer, implementation strategies, and common
failure modes.

The exposition is formal where formal algebra is the correct language and
explicit about the additional hypotheses needed when a transseries is to be
interpreted as an asymptotic expansion of an actual function.
\end{abstract}

\chapter*{Source and scope note}
\addcontentsline{toc}{chapter}{Source and scope note}
This monograph is an original synthesis based on the supplied archive, which
contains TeX and PDF versions of Gerald Edgar's foundational expositions on
transseries and composition, later structural papers on logarithmic and
surreal transseries, and several extensive tutorial drafts.  Additional
sources were checked online, especially Corless--Gonnet--Hare--Jeffrey--Knuth
on Lambert $\LambertW$, the NIST Digital Library of Mathematical Functions,
Salvy--Shackell on multiseries inversion, van der Hoeven's computational
transseries theory, and the power--log transseries literature associated with
Dulac series.  Full references appear in the bibliography.

The phrase \emph{polynomial--logarithmic transseries} is used here in a
specific computational sense.  The basic objects have the form
\[
 \sum_{n\ge n_0}X^{-n}p_n(\log X),
 \qquad p_n\in\K[L],
\]
possibly with fractional powers, Laurent-series coefficients in
$1/\log X$, or finitely many iterated logarithmic levels.  This is a useful
subclass of logarithmic--exponential transseries, not a proposed replacement
for the full theory.  It is large enough for Lambert $\LambertW$, many inverse
problems, quantile asymptotics, saddle-point calculations, and power--log
normal-form calculations, while remaining concrete enough for exact
coefficient algorithms.

\chapter*{How to read this article}
\addcontentsline{toc}{chapter}{How to read this article}
Readers mainly interested in calculation can proceed as follows.
\begin{enumerate}
\item Read Chapters~\ref{ch:scale}--\ref{ch:formal-ring} for the scale and
      notation.
\item Use Chapters~\ref{ch:multiplication} and~\ref{ch:division} for product,
      reciprocal, and quotient algorithms.
\item Use Chapters~\ref{ch:composition-generators} and
      \ref{ch:composition-taylor} for substitution.
\item Use Chapters~\ref{ch:reversion-fixed}--\ref{ch:reversion-lagrange} for
      compositional inverses.
\item Read Part~\ref{part:lambert} for the complete Lambert-$\LambertW$ calculation.
\item Keep Appendix~\ref{app:quick-reference} open as a formula sheet.
\end{enumerate}
The algebraic statements are given over a characteristic-zero coefficient
field $\K$, usually $\RealNumbers$ or $\ComplexNumbers$.  For an analytic interpretation, $X$ tends
to $+\infty$ unless a complex sector or another limit is explicitly stated.

\tableofcontents

\mainmatter

\part{The polynomial--logarithmic scale}

\chapter{Why powers and logarithms belong in one expansion}
\label{ch:scale}

\section{The prototype from Lambert's function}
Lambert's $\LambertW$ function is defined by
\begin{equation}\label{eq:W-def-intro}
 \LambertW(x)\EulerE^{\LambertW(x)}=x.
\end{equation}
Taking a logarithm gives
\begin{equation}\label{eq:W-log-intro}
 \LambertW(x)+\log \LambertW(x)=\log x.
\end{equation}
The dominant balance is $\LambertW(x)\sim\log x$, but substitution of this first
approximation leaves the smaller error $\log\log x$.  Thus the first two
scales are forced:
\[
 \LambertW(x)=\log x-\log\log x+\hbox{smaller terms}.
\]
If
\[
 L_1=\log x,\qquad L_2=\log L_1,
\]
then the complete expansion is organized in inverse powers of $L_1$, with a
polynomial in $L_2$ in each block:
\begin{equation}\label{eq:W-shape-intro}
 \LambertW(x)=L_1-L_2+\sum_{n\ge1}\frac{\LambertPolynomial{n}(L_2)}{L_1^n}.
\end{equation}
This block structure is the basic subject of the article.

\section{The hierarchy inside a block expansion}
Let $X\to+\infty$ and put
\[
 t=X^{-1},\qquad L=\log X.
\]
Then
\[
 t\precdom L^{-m}\precdom 1\precdom L^m\precdom t^{-1}
 \qquad(m>0).
\]
More generally,
\begin{equation}\label{eq:monomial-order}
 t^nL^k\succdom t^mL^j
 \quad\Longleftrightarrow\quad
 n<m\ \hbox{or}\ (n=m\ \hbox{and}\ k>j).
\end{equation}
The power of $t$ is compared first.  Only when the $t$ powers agree do we
compare logarithmic powers.  This is lexicographic asymptotic ordering.

For example,
\[
 \frac{L^{1000}}{X^3}\precdom\frac{1}{X^2L^{1000}},
\]
because one additional inverse power of $X$ defeats every fixed power of
$\log X$.  This elementary fact is the reason it is natural to store all
powers of $L$ at a fixed $X^{-n}$ as one coefficient polynomial.

\section{What ``transseries'' adds to an ordinary asymptotic series}
An ordinary inverse-power expansion
\[
 a_0+a_1X^{-1}+a_2X^{-2}+\cdots
\]
uses the scale $1,X^{-1},X^{-2},\ldots$.  A polynomial--logarithmic expansion
uses the refined scale
\[
 \cdots, X^{-n}L^2,\ X^{-n}L,\ X^{-n},\ X^{-n}L^{-1},\ldots,
\]
and may place another asymptotic series inside the coefficient of each
$X^{-n}$.  A full logarithmic--exponential transseries goes much further,
allowing exponentials, nested exponentials, arbitrary well-based supports,
and many levels of logarithms.  The present class is a computationally
transparent corner of that larger ordered field; see Edgar
\cite{EdgarBeginners,EdgarComposition} and van der Hoeven
\cite{vdHoeven2006} for the ambient theory.

\section{Relation to power--log and Dulac series}
In local dynamics one often studies $x\to0^+$ and writes
\[
 \ell(x)=-\frac{1}{\log x}.
\]
Power--log transseries then use monomials $x^\alpha\ell^k$.  Under $X=1/x$,
\[
 x^\alpha\ell^k=X^{-\alpha}(\log X)^{-k}.
\]
Thus the power--log notation and the present large-$X$ notation encode the
same two-scale geometry, with different choices of small generators.  Dulac
series and their extensions permit real exponent sets and logarithmic
Laurent coefficients; see Marde\v{s}i\'c--Resman--Rolin--\v{Z}upanovi\'c
\cite{MRRZ2016}.  Our integer-lattice class is deliberately narrower at
first, because this is where multiplication, division, composition, and
reversion can be displayed with exact finite algorithms.

\section{Formal equality, asymptotic equality, and analytic equality}
Three notions must be kept separate.
\begin{enumerate}
\item A \emph{formal identity} is an identity between coefficient arrays.
\item An \emph{asymptotic expansion} states that finite truncations approximate
      a function in a specified limit.
\item An \emph{analytic identity} is an equality of functions on a domain.
\end{enumerate}
Formal manipulation does not by itself prove that a series converges.  On the
other hand, convergence is not required for formal coefficient extraction.
Lambert $\LambertW$ is unusually friendly: its standard double logarithmic series is
not only asymptotic but is absolutely convergent in suitable real and complex
domains; see Kalugin and Jeffrey~\cite{KaluginJeffrey2012}, together with the
classical treatment of Corless et al.~\cite{Corless1996} and the branch data
collected in DLMF~\cite{DLMF413}.
Most of the algebra developed below does not rely on this special fact.

\begin{summarybox}
Polynomial--logarithmic transseries are block expansions in which inverse
powers of a large variable form the outer scale and powers or series in its
logarithm form inner coefficients.  The outer power is asymptotically decisive.
The same formalism appears at $0$ as power--log or Dulac-type series after a
change of variables.
\end{summarybox}

\chapter{The formal one-logarithm algebra}
\label{ch:formal-ring}

\section{The polynomial-coefficient Laurent algebra}
Let $\K$ be a field of characteristic zero.  Introduce two formal symbols
$t$ and $L$, to be interpreted later as $t=X^{-1}$ and $L=\log X$.

\begin{definition}[One-logarithm Laurent algebra]
The basic polynomial--logarithmic algebra is
\begin{equation}\label{eq:PL-poly-def}
 \PLpoly=\K[L]((t))
 =\SetOf{\sum_{n\ge n_0}p_n(L)t^n:
 n_0\in\IntegerNumbers,\ p_n\in\K[L]}.
\end{equation}
Only finitely many negative powers of $t$ are allowed, while infinitely many
nonnegative powers may occur.  Each coefficient $p_n$ is an ordinary finite
polynomial in $L$.
\end{definition}

The representation is unique.  It is useful to expand it further as
\[
 F=\sum_{n\ge n_0}\sum_{k=0}^{d_n}a_{n,k}t^nL^k.
\]
The support is the set of pairs $(n,k)$ for which $a_{n,k}\ne0$.
The asymptotic order is the lexicographic order in \eqref{eq:monomial-order}.

\begin{definition}[Leading data]
For nonzero $F\in\PLpoly$, let $n_0$ be the smallest $n$ with
$p_n\ne0$, and let $d=\deg p_{n_0}$.  Then
\[
 \lm(F)=t^{n_0}L^d,
 \qquad
 \lc(F)=\CoefficientExtraction{L}{d}p_{n_0}.
\]
The pair $(n_0,-d)$ is the full lexicographic valuation of $F$.
\end{definition}

\section{Why the support is legitimate}
At fixed $n$, only finitely many $k$ occur.  The outer exponents are bounded
below.  Consequently every nonempty support has a largest monomial in the
asymptotic order, and each coefficient of a product receives only finitely
many contributions.  This is the elementary lattice version of the
Hahn--Neumann support mechanism underlying general transseries.

\begin{proposition}[Leading-term arithmetic]\label{prop:leading-arithmetic}
For nonzero $F,G\in\PLpoly$,
\[
 \lm(FG)=\lm(F)\lm(G),
 \qquad
 \lc(FG)=\lc(F)\lc(G).
\]
If $\lm(F)\ne\lm(G)$, then the larger of the two leading monomials is the
leading monomial of $F+G$.
\end{proposition}
\begin{proof}
Write the leading terms as $a t^nL^k$ and $b t^mL^j$.  Their product
$ab t^{n+m}L^{k+j}$ is larger than every product involving a lower monomial
of either factor.  The assertion for sums is immediate unless the leading
monomials agree, in which case cancellation may occur and one inspects the
next term.
\end{proof}

\section{Block order and truncation}
For coefficient algorithms we use the outer $t$-adic filtration.

\begin{definition}[Block order]
For $F=\sum_{n\ge n_0}p_n(L)t^n\ne0$, define
\[
 \nu_tF=n_0.
\]
Write $F=\FormalBigOAt{t^N}{t}$ when $\nu_tF\ge N$.  This is a formal block statement,
not the analytic estimate
$\AbsoluteValue{F}\le C\AbsoluteValue{t}^N$.
\end{definition}

The truncation through block $N$ is
\[
 \trunc{F}{N}=\sum_{n=n_0}^{N}p_n(L)t^n,
\]
and
\[
 F=G+\FormalBigOAt{t^{N+1}}{t}
 \quad\Longleftrightarrow\quad
 \trunc{F}{N}=\trunc{G}{N}.
\]
All finite-order algorithms below retain a fixed finite principal part and identify
two tails when their difference lies in
\(t^{N+1}\K[L][[t]]\).  This is a valuation-filtration equivalence on that
fixed-principal-part slice, not a quotient algebra of \(\K[L]((t))\): the Laurent
variable \(t\) is a unit, so the ideal generated by \(t^{N+1}\) in the full Laurent
ring is the whole ring.

\section{Why a field completion is useful}
The ring $\K[L]$ has only constant units.  Thus $L+1$ is not invertible in
$\K[L]$, even though as an asymptotic object
\[
 \frac{1}{L+1}=L^{-1}-L^{-2}+L^{-3}-\cdots.
\]
Exact rational dependence can first be retained in the rational-log Laurent
field
\[
 \RationalLogLaurentField{\K}=\K(L)((t)).
\]
Expanding each rational coefficient at \(L=\infty\) then embeds this field into
the distinct Laurent-completed coefficient field in \(L^{-1}\):
\begin{equation}\label{eq:PL-field-def}
 \mathcal C_L=\K((L^{-1})),
 \qquad
 \PLfield=\mathcal C_L((t)).
\end{equation}
An element of $\mathcal C_L$ may contain finitely many positive powers of $L$
and an infinite tail of negative powers.  Therefore
\[
 \K[L]\subset\K[L,L^{-1}]\subset\K((L^{-1})).
\]
The outer field $\PLfield$ is an iterated Laurent-series field and hence a
field.  The embedding replaces, for example, \((L+1)^{-1}\) by its displayed
infinite Laurent tail; it is not an identification of the two coefficient
models, and a finite computation in \(\PLfield\) needs an inner
\(L^{-1}\)-cutoff in addition to its outer \(t\)-cutoff.

\begin{remark}
The order of iteration matters.  In $\K((L^{-1}))((t))$, the generator $t$ is
smaller than every power of $L^{-1}$.  This matches the asymptotic fact that
$X^{-1}=\LittleO((\log X)^{-m})$ for every fixed $m$.
\end{remark}

\section{Fractional powers and general exponent lattices}
If a problem contains $X^{-1/q}$, one may replace $t$ by $s=X^{-1/q}$ and use
$\K((L^{-1}))((s))$.  More generally, for an additive subgroup
$\Gamma\subset\RealNumbers$ one may consider
\[
 \sum_{\gamma\in S}p_\gamma(L)X^{-\gamma},
\]
where $S\subset\Gamma$ is well ordered and satisfies a finite-decomposition
condition.  Multiplication is then a Hahn-series convolution.  Grid-based
transseries place $S$ inside a finitely generated additive semigroup, which is
especially convenient for computation \cite{EdgarBeginners,EdgarComposition}.
The integer-lattice notation will be used in the main derivations, since the
same formulas then extend verbatim after replacing integer indices by
admissible exponents.

\section{Nested logarithmic levels}
Let
\[
 L_0=X,\qquad L_{j+1}=\log L_j,
 \qquad \InverseLogScale{j}=L_j^{-1}.
\]
For fixed depth $d$, a natural iterated field is
\begin{equation}\label{eq:nested-field}
 \K((\InverseLogScale{d}))\cdots
 ((\InverseLogScale{2}))((\InverseLogScale{1})).
\end{equation}
The rightmost variable $\InverseLogScale{1}$ is the outermost and smallest
scale.  In
particular,
\[
 \InverseLogScale{1}\precdom \InverseLogScale{2}^m
 \precdom\InverseLogScale{3}^n\precdom1
 \qquad(m,n>0).
\]
The Lambert expansion \eqref{eq:W-shape-intro} belongs to the subring
\[
 \K[L_2]((L_1^{-1}))
\]
with base variable $X=L_1$ and logarithm $L=L_2$.

\section{Asymptotic interpretation and uniqueness}
For actual functions, a convenient one-layer definition is the following.

\begin{definition}[Block asymptotic expansion]
We write
\begin{equation}\label{eq:block-asymptotic}
 f(X)\PolynomialLogAsymptotic
 \sum_{n\ge n_0}p_n(\log X)X^{-n}
\end{equation}
if for every integer $N\ge n_0$,
\[
 f(X)-\sum_{n=n_0}^{N}p_n(\log X)X^{-n}=\LittleO(X^{-N})
 \qquad(X\to+\infty).
\]
\end{definition}

\begin{proposition}[Uniqueness of polynomial blocks]
An expansion in the sense of \eqref{eq:block-asymptotic} has unique
coefficient polynomials $p_n$.
\end{proposition}
\begin{proof}
Suppose two expansions first differ at block $n$.  Their difference divided by
$X^{-n}$ is a nonzero polynomial $Q(\log X)+\LittleO(1)$.  A nonzero polynomial in
$\log X$ cannot tend to zero as $X\to\infty$, a contradiction.
\end{proof}

For coefficients in $\K((L^{-1}))$, one uses nested truncations: first truncate
in $t$, then within each retained coefficient truncate in $L^{-1}$.  This is
the analytic counterpart of the iterated field \eqref{eq:PL-field-def}.

\begin{summarybox}
The elementary polynomial--logarithmic ring is $\K[L]((t))$ with
$t=X^{-1}$ and $L=\log X$.  It is an ordered integral domain and a complete
$t$-adic algebra, but not a field.  The iterated Laurent field
$\K((L^{-1}))((t))$ supplies the missing inverses while preserving the rule
that $t$ is smaller than every power of $L^{-1}$.
\end{summarybox}

\part{Arithmetic and differential calculus}

\chapter{Multiplication: convolution by logarithmic blocks}
\label{ch:multiplication}

\section{The block Cauchy product}
Let
\begin{equation}\label{eq:FG-blocks}
 F=\sum_{n\ge n_0}A_n(L)t^n,
 \qquad
 G=\sum_{m\ge m_0}B_m(L)t^m.
\end{equation}
Their product is obtained by the ordinary Cauchy rule in the outer variable:
\begin{equation}\label{eq:block-convolution}
 FG=\sum_{N\ge n_0+m_0}C_N(L)t^N,
 \qquad
 C_N(L)=\sum_{n+m=N}A_n(L)B_m(L).
\end{equation}
For a fixed $N$, the restrictions $n\ge n_0$ and $m\ge m_0$ leave only
finitely many pairs.  Hence \eqref{eq:block-convolution} is an exact algebraic
formula, not a rearrangement of an analytically convergent double series.

\begin{theorem}[Closure and finite determination]\label{thm:mult-closure}
The algebra $\PLpoly=\K[L]((t))$ is closed under multiplication.  Moreover,
the coefficient of $t^N$ in $FG$ depends only on the blocks
\[
 A_{n_0},\ldots,A_{N-m_0},
 \qquad
 B_{m_0},\ldots,B_{N-n_0}.
\]
Consequently, if
$F=\ApproximationOf{F}+\FormalBigOAt{t^{p}}{t}$ and
$G=\ApproximationOf{G}+\FormalBigOAt{t^{q}}{t}$, then
\begin{equation}\label{eq:mult-error-rule}
 FG-\ApproximationOf{F}\ApproximationOf{G}
 =\FormalBigOAt{t^{\min(p+m_0,q+n_0)}}{t}.
\end{equation}
\end{theorem}
\begin{proof}
Formula \eqref{eq:block-convolution} is a finite sum of products of
polynomials and therefore gives $C_N\in\K[L]$.  The support remains bounded
below by $n_0+m_0$.  For the error rule, write
$F=\ApproximationOf{F}+E$ and $G=\ApproximationOf{G}+H$, with
$\nu_tE\ge p$ and $\nu_tH\ge q$.  Then
\[
 FG-\ApproximationOf{F}\ApproximationOf{G}=E\ApproximationOf{G}+\ApproximationOf{F} H+EH,
\]
and the three outer orders are at least $p+m_0$, $q+n_0$, and $p+q$.
The first two bounds give \eqref{eq:mult-error-rule}.
\end{proof}

\section{The coefficient-level convolution}
Write
\[
 A_n(L)=\sum_{i=0}^{a_n}a_{n,i}L^i,
 \qquad
 B_m(L)=\sum_{j=0}^{b_m}b_{m,j}L^j.
\]
Then
\begin{equation}\label{eq:two-dimensional-convolution}
 [t^NL^k](FG)
 =\sum_{n+m=N}\ \sum_{i+j=k}a_{n,i}b_{m,j}.
\end{equation}
Thus multiplication is a two-dimensional convolution, but it is normally
more efficient to perform it in two levels: first convolve the outer blocks,
then multiply the coefficient polynomials.  This respects the asymptotic
hierarchy and allows a separate cutoff in $L^{-1}$ if the coefficient field is
$\K((L^{-1}))$.

\begin{proposition}[Degree and support bounds]\label{prop:degree-bounds}
For the product coefficient in \eqref{eq:block-convolution},
\[
 \degL C_N\le
 \max_{n+m=N}\bigl(\degL A_n+\degL B_m\bigr).
\]
If all $A_n$ and $B_m$ have nonnegative coefficients over an ordered field,
then equality holds whenever the maximizing blocks are nonzero.
\end{proposition}
\begin{proof}
Each summand has degree at most the displayed sum of degrees.  Cancellation
can lower the degree in general.  With nonnegative coefficients there is no
cancellation at the largest power of $L$.
\end{proof}

\section{A complete multiplication example}
Take
\begin{align*}
 F&=1+(L+1)t+(L^2-L)t^2+(2L+3)t^3+\FormalBigOAt{t^4}{t},\\
 G&=1-Lt+(2L-1)t^2+(L^2+1)t^3+\FormalBigOAt{t^4}{t}.
\end{align*}
The blocks are computed one at a time:
\begin{align*}
 C_0&=1,\\
 C_1&=(L+1)-L=1,\\
 C_2&=(L^2-L)+(L+1)(-L)+(2L-1)
      =L-1,\\
 C_3&=(2L+3)+(L^2-L)(-L)+(L+1)(2L-1)+(L^2+1)\\
    &=-L^3+4L^2+2L+3.
\end{align*}
Therefore
\begin{equation}\label{eq:product-example}
 FG=1+t+(L-1)t^2+(-L^3+4L^2+2L+3)t^3+\FormalBigOAt{t^4}{t}.
\end{equation}
The calculation illustrates why a block representation is preferable to a
flat list of monomials: the outer convolution is transparent, and all
logarithmic algebra is delegated to ordinary polynomial arithmetic.

\section{Truncated multiplication algorithm}
Suppose both factors begin at block zero and are wanted through $t^N$.

\begin{algorithm}[H]
\caption{Truncated polynomial--logarithmic product}
\label{alg:pl-product}
\KwIn{Blocks $A_0,\ldots,A_N$ and $B_0,\ldots,B_N$ in a coefficient ring}
\KwOut{Blocks $C_0,\ldots,C_N$ of $\trunc{FG}{N}$}
\For{$r\leftarrow0$ \KwTo $N$}{
  $C_r\leftarrow0$\;
  \For{$i\leftarrow0$ \KwTo $r$}{
    $C_r\leftarrow C_r+A_iB_{r-i}$\;
  }
}
\Return{$(C_0,\ldots,C_N)$}\;
\end{algorithm}

With schoolbook polynomial multiplication, the arithmetic cost is governed by
both the number of blocks and their logarithmic degrees.  If the typical
coefficient degree is $d$, the direct cost is $\BigO(N^2d^2)$ scalar operations.
Fast convolution and fast polynomial multiplication can reduce this, but for
asymptotic calculations the degrees and truncation orders are often modest,
and sparse storage is more important than asymptotic complexity.

\section{Multiplication in the completed coefficient field}
Let
\[
 A_n(L)=\sum_{j\ge j_n}a_{n,j}L^{-j},
 \qquad
 B_m(L)=\sum_{k\ge k_m}b_{m,k}L^{-k}.
\]
At a fixed outer block $N$, the product coefficient $C_N$ is a finite sum of
Laurent-series products.  At a fixed inner power $L^{-r}$, each such product
has only finitely many decompositions $j+k=r$.  Thus the iterated order
``first $t$, then $L^{-1}$'' makes multiplication coefficientwise finite at
both levels.

\begin{warning}
A rectangular cutoff in the exponent plane need not respect asymptotic order.
For example, discarding all terms with $\AbsoluteValue{k}>10$ may remove
$t^2L^{100}$ while
retaining the much smaller $t^3L^{-100}$.  The safe strategy is a block cutoff
in $t$, followed by a separately declared cutoff inside each coefficient.
\end{warning}

\section{Products of asymptotic expansions}
Formal multiplication mirrors a standard analytic rule.

\begin{proposition}[Analytic product rule]\label{prop:analytic-product}
Assume
\[
 f(X)=\sum_{n=n_0}^{N}A_n(L)X^{-n}+\LittleO(X^{-N}),
 \quad
 g(X)=\sum_{m=m_0}^{M}B_m(L)X^{-m}+\LittleO(X^{-M}),
\]
where the displayed finite sums have at most polynomial growth in $L$.  Then
for every $R$ below the available precision,
\[
 f(X)g(X)=\sum_{n+m\le R}A_n(L)B_m(L)X^{-n-m}
 +\LittleO(X^{-R}).
\]
\end{proposition}
\begin{proof}
Multiply the two finite expansions and collect blocks.  A discarded term
$X^{-p}$ times a polynomial in $L$ is $\LittleO(X^{-p+\delta})$ for every
$\delta>0$.  Taking the input expansions sufficiently far beyond the desired
integer block $R$ absorbs all logarithmic factors.  The stated block
hypothesis, which uses exact integer powers, then yields the result.
\end{proof}

\begin{summarybox}
Multiplication is ordinary Cauchy convolution in $t=X^{-1}$ with polynomial or
Laurent-series arithmetic in $L=\log X$.  Every retained block receives only
finitely many contributions.  Truncation should always follow the iterated
asymptotic hierarchy rather than an arbitrary rectangular exponent cutoff.
\end{summarybox}


\chapter{Reciprocals and division}
\label{ch:division}

\section{Which elements are units?}
A nonzero element of the field $\PLfield=\K((L^{-1}))((t))$ is invertible.
In the smaller ring $\PLpoly=\K[L]((t))$, invertibility is more restrictive.

\begin{proposition}[Units of the polynomial coefficient ring]
Let
\[
 F=t^\nu\bigl(A_0(L)+A_1(L)t+A_2(L)t^2+\cdots\bigr),
 \qquad A_0\ne0.
\]
Then $F$ is a unit of $\K[L]((t))$ if and only if $A_0$ is a nonzero constant.
It is always a unit of $\K((L^{-1}))((t))$.
\end{proposition}
\begin{proof}
The factor $t^\nu$ is invertible.  In a formal power-series ring $R[[t]]$, a
series is a unit precisely when its constant coefficient is a unit of $R$.
The only units of $\K[L]$ are nonzero constants, whereas every nonzero element
of the field $\K((L^{-1}))$ is a unit.
\end{proof}

Thus division may enlarge the coefficient class.  For example,
\[
 \frac1{L+1}=L^{-1}-L^{-2}+L^{-3}-\cdots
\]
is perfectly meaningful in $\K((L^{-1}))$, but not in $\K[L]$.

\section{Normalization and the geometric expansion}
Every nonzero $B\in\PLfield$ can be written uniquely as
\begin{equation}\label{eq:division-normalization}
 B=b\,t^\nu L^\mu(1+S),
 \qquad b\in\K^\times,
 \qquad S\precdom1.
\end{equation}
Then
\begin{equation}\label{eq:inverse-geometric}
 B^{-1}=b^{-1}t^{-\nu}L^{-\mu}
 \sum_{r\ge0}(-S)^r.
\end{equation}
The support conditions guarantee that the geometric family is summable: for
each retained monomial, only finitely many powers of $S$ can contribute.
Formula \eqref{eq:inverse-geometric} is the universal conceptual construction.
For finite computation, a coefficient recurrence is usually faster.

\section{The triangular quotient recurrence}
First assume
\begin{equation}\label{eq:quotient-normalized}
 A=\sum_{n\ge0}a_n t^n,
 \qquad
 B=\sum_{n\ge0}b_n t^n,
 \qquad b_0\ne0,
\end{equation}
with coefficients in a field such as $\K((L^{-1}))$.  Write
$Q=A/B=\sum_{n\ge0}q_nt^n$.  The identity $BQ=A$ gives
\begin{equation}\label{eq:quotient-recurrence}
 q_0=\frac{a_0}{b_0},
 \qquad
 q_N=\frac{1}{b_0}
 \left(a_N-\sum_{j=1}^{N}b_jq_{N-j}\right),
 \quad N\ge1.
\end{equation}
This is triangular: once $q_0,\ldots,q_{N-1}$ are known, $q_N$ follows by one
coefficient-field division.

For a reciprocal $R=B^{-1}=\sum r_nt^n$, set $a_0=1$ and $a_N=0$ for
$N\ge1$:
\begin{equation}\label{eq:reciprocal-recurrence}
 r_0=b_0^{-1},
 \qquad
 r_N=-b_0^{-1}\sum_{j=1}^{N}b_jr_{N-j}.
\end{equation}

\begin{theorem}[Correctness of truncated division]\label{thm:division-correct}
If $q_0,\ldots,q_N$ are defined by \eqref{eq:quotient-recurrence}, then
\[
 B\sum_{n=0}^{N}q_nt^n
 =A+\FormalBigOAt{t^{N+1}}{t}.
\]
They are the unique blocks with this property.
\end{theorem}
\begin{proof}
The coefficient of $t^N$ in $BQ$ is
$b_0q_N+\sum_{j=1}^Nb_jq_{N-j}$, which equals $a_N$ by construction.
Induction gives all blocks through $N$.  If two truncated quotients worked,
their difference would have a first nonzero block $ct^r$; multiplication by
$B$ would have first block $b_0ct^r\ne0$, a contradiction.
\end{proof}

\section{A reciprocal example}
Let
\[
 B=1-Lt+(2L-1)t^2+(L^2+1)t^3+\FormalBigOAt{t^4}{t}.
\]
Applying \eqref{eq:reciprocal-recurrence} gives
\begin{align*}
 r_0&=1,\\
 r_1&=L,\\
 r_2&=L^2-2L+1,\\
 r_3&=L^3-5L^2+2L-1,\\
 r_4&=L^4-8L^3+7L^2-6L+1,
\end{align*}
where $r_4$ is determined by taking the omitted $t^4$ block of $B$ to be
zero.  Hence
\begin{align}\label{eq:reciprocal-example}
 B^{-1}={}&1+Lt+(L^2-2L+1)t^2
 +(L^3-5L^2+2L-1)t^3\notag\\
 &+(L^4-8L^3+7L^2-6L+1)t^4+\FormalBigOAt{t^5}{t}.
\end{align}
Multiplication by $B$ checks the result block by block.

\section{A quotient example}
With
\begin{align*}
 A&=1+(L+1)t+(L^2-L)t^2+(2L+3)t^3+\FormalBigOAt{t^4}{t},\\
 B&=1-Lt+(2L-1)t^2+(L^2+1)t^3+\FormalBigOAt{t^4}{t},
\end{align*}
recurrence \eqref{eq:quotient-recurrence} yields
\begin{align}\label{eq:quotient-example}
 \frac AB={}&1+(2L+1)t+(3L^2-2L+1)t^2\notag\\
 &+(3L^3-7L^2+3L+3)t^3+\FormalBigOAt{t^4}{t}.
\end{align}
A reliable implementation should verify the quotient by multiplying the
right-hand side by the divisor to the requested order.  This residual check
is inexpensive and catches most cutoff or normalization mistakes.

\section{Leading monomials other than one}
For general $A$ and $B$, first factor the leading monomial of $B$:
\[
 B=b t^\nu L^\mu(1+S).
\]
Then
\[
 \frac AB=b^{-1}t^{-\nu}L^{-\mu}A(1+S)^{-1}.
\]
If the leading coefficient is a nonmonomial Laurent series $b_0(L)$, invert
it in $\K((L^{-1}))$ before applying the outer recurrence.  For example,
\begin{align*}
 \frac{1}{L^2+3L+2}
 &=L^{-2}\frac1{1+3L^{-1}+2L^{-2}}\\
 &=L^{-2}-3L^{-3}+7L^{-4}-15L^{-5}+31L^{-6}+\cdots.
\end{align*}
There are therefore two nested divisions: an inner Laurent-series inversion
for the leading coefficient and an outer $t$-series inversion.

\section{Newton inversion}
For high order, Newton iteration replaces the linear recurrence by a
precision-doubling scheme.  If $R$ approximates $B^{-1}$, set
\begin{equation}\label{eq:newton-reciprocal}
 R_{\mathrm{new}}=R(2-BR).
\end{equation}
If $BR=1-E$, then
\[
 BR_{\mathrm{new}}=(1-E)(1+E)=1-E^2.
\]
Thus an error of outer order $m$ becomes one of order $2m$.  Starting with
$R_0=b_0^{-1}$, one truncates after each step to orders
$1,2,4,8,\ldots$ until the target precision is reached.

\begin{algorithmicprinciple}[Guard precision]
When a quotient will later be multiplied by a factor of negative outer order,
compute it beyond the final requested block.  Truncation is safe only after
propagating valuations through the entire expression tree.  A robust system
uses explicit precision objects rather than a single global cutoff.
\end{algorithmicprinciple}

\section{Division of asymptotic expansions}
Suppose $g(X)$ has a nonzero leading block and is bounded away from zero after
factoring its leading monomial.  If formal division gives
$Q_N(X)$ through block $N$ and
\[
 f(X)-g(X)Q_N(X)=\LittleO(X^{-N}),
\]
then
\[
 \frac{f(X)}{g(X)}-Q_N(X)
 =\frac{\LittleO(X^{-N})}{g(X)}.
\]
The leading scale of $g$ converts this residual into the corresponding
quotient error.  This simple residual argument is often more reliable than
trying to manipulate analytic remainder symbols at every intermediate step.

\begin{summarybox}
Normalize a divisor by its leading monomial and invert the remaining unit.
At finite order, quotient blocks satisfy a triangular recurrence.  Newton's
formula $R\mapsto R(2-BR)$ doubles precision and is preferable at high order.
Division may force polynomial logarithmic coefficients to expand into Laurent
series in $1/\log X$.
\end{summarybox}


\chapter{Powers, logarithms, exponentials, and closure boundaries}
\label{ch:powers}

\section{General powers of a normalized element}
Let
\begin{equation}\label{eq:power-normalized}
 F=c\,t^\nu L^\mu(1+S),
 \qquad S\precdom1.
\end{equation}
For $\alpha\in\K$ one formally defines
\begin{equation}\label{eq:general-power}
 F^\alpha=c^\alpha t^{\alpha\nu}L^{\alpha\mu}
 \sum_{r\ge0}\binom{\alpha}{r}S^r.
\end{equation}
Over $\RealNumbers$, a real value for $c^\alpha$ requires the usual sign conditions.
Over $\ComplexNumbers$, an analytic logarithm branch $\mathcal B_c$ must be
selected and
$c^\alpha=\exp\!\bigl(\alpha\LogarithmOnBranch{\mathcal B_c}{c}\bigr)$.
Rational $\alpha$ may
force a fractional outer exponent lattice; irrational $\alpha$ requires a
more general ordered exponent group.

\begin{proposition}[Finite determination of a binomial power]
If $S=\FormalBigOAt{t^q}{t}$ with $q>0$, then the coefficient of $t^N$ in
$(1+S)^\alpha$ receives contributions only from powers
$r\le\Floor{N/q}$.  Hence \eqref{eq:general-power} is
coefficientwise finite.
\end{proposition}
\begin{proof}
The outer order of $S^r$ is at least $rq$.
\end{proof}

\section{Logarithms}
For \eqref{eq:power-normalized},
\begin{equation}\label{eq:log-normalized}
 \log F=\log c+\nu\log t+\mu\log L+\log(1+S).
\end{equation}
Since $\log t=-L$, this becomes
\begin{equation}\label{eq:log-normalized-expanded}
 \log F=\log c-\nu L+\mu\log L
 +\sum_{r\ge1}\frac{(-1)^{r+1}}rS^r.
\end{equation}
Two closure facts are visible.
\begin{enumerate}
\item If $\mu=0$ and $S$ is $t$-adically small, the logarithm stays in the
      one-logarithm class, apart from the allowed linear term in $L$.
\item If $\mu\ne0$, the new generator $\log L=\log\log X$ appears.  The
      one-logarithm class is therefore not closed under logarithm, but the
      finite nested-logarithm tower is.
\end{enumerate}

\begin{example}
For $F=X^3L^{-2}(1+(L+1)/X)$,
\begin{align*}
 \log F
 &=3L-2\log L+\log\left(1+(L+1)t\right)\\
 &=3L-2\log L+(L+1)t-\frac12(L+1)^2t^2
   +\frac13(L+1)^3t^3+\cdots.
\end{align*}
The nested logarithm is forced by the factor $L^{-2}$, not by the small
correction.
\end{example}

\section{Exponentials}
If $S\precdom1$, then
\begin{equation}\label{eq:exp-small}
 \exp S=\sum_{r\ge0}\frac{S^r}{r!}
\end{equation}
is summable and begins with $1$.  If
$F=F_{\mathrm{large}}+c+S$ is decomposed into a purely large part, constant,
and small part, then
\[
 \exp F=\EulerE^c\exp(F_{\mathrm{large}})\exp S.
\]
The factor $\exp(F_{\mathrm{large}})$ is generally a new transmonomial.
For example,
\[
 \exp(L)=X=t^{-1}
\]
stays inside the class, but $\exp(L^2)=\exp((\log X)^2)$ does not belong to
$\K((L^{-1}))((t))$.  It belongs naturally to a larger
logarithmic--exponential transseries field.

\begin{pitfallbox}
The statement ``$S$ is small'' must refer to the asymptotic ordering, not to
small numerical coefficients.  The formal exponential of $L^{-1}$ is a unit,
whereas the exponential of $L^2$ creates a new dominant scale.  Treating both
with the same flat power-series routine destroys the monomial hierarchy.
\end{pitfallbox}

\section{Analytic functions of a unit}
Let $\Phi(z)=\sum_{r\ge0}c_r(z-z_0)^r$ be analytic or merely formal at
$z_0$, and let $F=z_0+S$ with $S\precdom1$.  Then
\begin{equation}\label{eq:analytic-substitution-unit}
 \Phi(F)=\sum_{r\ge0}c_rS^r
\end{equation}
is a valid transseries substitution whenever the powers are summable.  This
single principle supplies square roots, trigonometric functions of small
arguments, rational powers, and local inverse functions.

For example,
\begin{align*}
 (1+S)^{1/2}&=1+\frac12S-\frac18S^2+\frac1{16}S^3-\cdots,\\
 \frac{1}{1+S}&=1-S+S^2-S^3+\cdots,\\
 \log(1+S)&=S-\frac12S^2+\frac13S^3-\cdots.
\end{align*}
All three use the same support argument.

\section{A closure table}
\begin{center}
\begin{tabularx}{0.96\textwidth}{@{}lX X@{}}
\toprule
Operation & Typical hypothesis & Resulting class \\
\midrule
Addition, multiplication & none beyond support conditions
 & $\K[L]((t))$ \\
Reciprocal & leading coefficient constant
 & $\K[L]((t))$ \\
General reciprocal & nonzero leading coefficient
 & $\K((L^{-1}))((t))$ \\
Derivative & any polynomial--log series
 & same one-logarithm class \\
$\Phi(c+S)$ & $S\precdom1$, $\Phi$ local formal series
 & same class if coefficients permit \\
Logarithm & normalized $ct^\nu L^\mu(1+S)$
 & adds $\log L$ when $\mu\ne0$ \\
Exponential & small argument
 & a unit in the same class \\
Exponential & nontrivial large part
 & may require new exponential monomials \\
Composition & large admissible inner argument
 & may require fractional powers or deeper logs \\
\bottomrule
\end{tabularx}
\end{center}

\begin{summarybox}
Normalize before applying nonlinear operations.  Binomial, logarithmic, and
exponential series are safe on asymptotically small units.  Logarithms of
logarithmic powers create deeper logarithms, and exponentials of genuinely
large polynomial--logarithmic expressions usually leave the class.  These are
structural closure boundaries, not implementation defects.
\end{summarybox}


\chapter{Differentiation and antiderivatives}
\label{ch:differentiation}

\section{The basic derivation}
Recall $t=X^{-1}$ and $L=\log X$.  Therefore
\[
 \frac{\Differential t}{\Differential X}=-t^2,
 \qquad
 \frac{\Differential L}{\Differential X}=t.
\]
On a coefficient block this gives the fundamental formula
\begin{equation}\label{eq:derivative-block}
 \D\bigl(P(L)t^n\bigr)
 \DefinitionEquals\frac{\Differential}{\Differential X}\bigl(P(L)X^{-n}\bigr)
 =t^{n+1}\bigl(P'(L)-nP(L)\bigr).
\end{equation}
Equivalently,
\begin{equation}\label{eq:derivation-operator}
 \D=-t^2\frac{\partial}{\partial t}
     +t\frac{\partial}{\partial L}.
\end{equation}

\begin{proposition}[Closure under differentiation]
The derivation $\D$ maps $\K[L]((t))$ into itself and raises outer order by at
least one unless cancellation occurs in a leading coefficient.  It extends
termwise to $\K((L^{-1}))((t))$.
\end{proposition}
\begin{proof}
Formula \eqref{eq:derivative-block} maps a polynomial to a polynomial and
sends block $n$ to block $n+1$.  For Laurent-series coefficients,
differentiation with respect to $L$ also preserves $\K((L^{-1}))$.
\end{proof}

\section{The Euler derivation}
It is often cleaner to use
\begin{equation}\label{eq:euler-derivation}
 \EulerDerivation=X\frac{\Differential}{\Differential X}=t^{-1}\D.
\end{equation}
Then
\begin{equation}\label{eq:euler-block}
 \EulerDerivation\bigl(P(L)t^n\bigr)=t^n(P'-nP).
\end{equation}
Thus each outer block is invariant under $\EulerDerivation$, and its coefficient is
acted on by the first-order operator $\partial_L-n$.  This diagonal block
structure is useful in differential equations.

\section{Repeated derivatives}
Iterating \eqref{eq:derivative-block} gives a compact formula.

\begin{theorem}[Repeated derivative formula]\label{thm:repeated-derivative}
For $r\ge0$,
\begin{equation}\label{eq:repeated-derivative}
 \D^r\bigl(P(L)t^n\bigr)
 =t^{n+r}
 \prod_{j=0}^{r-1}\bigl(\partial_L-(n+j)\bigr)P(L),
\end{equation}
where the empty product is the identity.
\end{theorem}
\begin{proof}
The case $r=0$ is trivial.  If the formula holds for $r$, apply
\eqref{eq:derivative-block} to its coefficient at outer exponent $n+r$.
This appends the commuting factor $\partial_L-(n+r)$.
\end{proof}

In particular,
\begin{align*}
 \D^r(L^m)
 &=t^r\prod_{j=0}^{r-1}(\partial_L-j)L^m,\\
 \D^r(t^n)
 &=(-1)^r(n)^{\overline r}t^{n+r},
\end{align*}
where $(n)^{\overline r}=n(n+1)\cdots(n+r-1)$ is the rising factorial.
The first identity is the source of the Stirling-number polynomials in the
Lambert expansion.

\section{Product and chain rules}
Because \eqref{eq:derivation-operator} is an ordinary derivation on the
formal generators, it satisfies
\[
 \D(FG)=(\D F)G+F(\D G).
\]
Whenever composition is admissible,
\begin{equation}\label{eq:formal-chain-rule}
 \D(F\compose G)=(\D F\compose G)\D G.
\end{equation}
The chain rule can be verified first for the generators $t$ and $L$, then for
finite algebraic expressions, and finally extended by coefficientwise
summability.  In a full transseries field it is one of the central
compatibility properties of composition \cite{EdgarComposition,vdHoeven2006}.

\section{Solving for an antiderivative block}
To integrate $Q(L)t^m$ with $m\ne1$, seek an antiderivative of the form
$P(L)t^{m-1}$.  Equation \eqref{eq:derivative-block} requires
\begin{equation}\label{eq:antiderivative-coefficient}
 P'-(m-1)P=Q.
\end{equation}
If $Q$ is a polynomial and $m\ne1$, this equation has a unique polynomial
solution.  Indeed,
\begin{equation}\label{eq:antiderivative-finite}
 P=-\sum_{j=0}^{\deg Q}
 \frac{Q^{(j)}}{(m-1)^{j+1}}.
\end{equation}
Differentiating the finite sum verifies \eqref{eq:antiderivative-coefficient}.
Therefore
\begin{equation}\label{eq:integral-block}
 \int Q(L)X^{-m}\Differential X
 =X^{1-m}
 \left[-\sum_{j=0}^{\deg Q}\frac{Q^{(j)}(L)}{(m-1)^{j+1}}\right]+C,
 \qquad m\ne1.
\end{equation}

When $m=1$, equation \eqref{eq:antiderivative-coefficient} becomes $P'=Q$,
so
\[
 \int Q(L)X^{-1}\Differential X=\int Q(L)\Differential L,
\]
again polynomial.  Hence the polynomial--logarithmic ring is closed under
termwise antiderivatives, up to an arbitrary constant, in this one-level
integer-exponent setting.

\begin{example}
Since
\[
 \int \frac{L^2}{X^3}\Differential X=X^{-2}P(L),
\]
we solve $P'-2P=L^2$ and obtain
\[
 P=-\frac12L^2-\frac12L-\frac14.
\]
Thus
\[
 \int \frac{(\log X)^2}{X^3}\Differential X
 =-\frac{2L^2+2L+1}{4X^2}+C.
\]
\end{example}

\section{Differential equations by block matching}
Consider
\begin{equation}\label{eq:simple-ode}
 Y'+Y=\frac{L}{X}.
\end{equation}
Seek $Y=\sum_{n\ge1}p_n(L)t^n$.  From
\eqref{eq:derivative-block},
\[
 Y'+Y=p_1t+\sum_{n\ge2}
 \bigl(p_n+p_{n-1}'-(n-1)p_{n-1}\bigr)t^n.
\]
Hence $p_1=L$ and
\[
 p_n=(n-1)p_{n-1}-p_{n-1}',\qquad n\ge2.
\]
The first blocks are
\[
 Y=Lt+(L-1)t^2+(2L-3)t^3+(6L-11)t^4+\cdots.
\]
This example shows how differentiation's shift by one outer block makes many
linear differential equations triangular.

\section{Why differentiation controls composition and inversion}
If $A=\FormalBigOAt{1}{t}$ contains no $X$ term, then
\begin{equation}\label{eq:A-derivative-small}
 \D A=\FormalBigOAt{t}{t}.
\end{equation}
Consequently, changing an argument by an error $E=\FormalBigOAt{t^r}{t}$ changes
$A$ by roughly $A'E=\FormalBigOAt{t^{r+1}}{t}$.  This one-block gain is the mechanism
behind the contraction used for compositional inversion.  Repeated derivatives
also show that the $n$th Lagrange--B\"urmann term contributes no earlier than
block $n-1$.

\begin{summarybox}
The derivative of $P(L)X^{-n}$ is
$X^{-n-1}(P'-nP)$.  Repeated derivatives factor into commuting operators
$\partial_L-(n+j)$.  Differentiation raises outer order, a fact that makes
near-identity Taylor composition and Lagrange reversion coefficientwise
finite.  Polynomial--logarithmic blocks also admit explicit polynomial
antiderivatives.
\end{summarybox}

\part{Composition}

\chapter{Composition by substitution of the generators}
\label{ch:composition-generators}

\section{Why composition is the delicate operation}
Addition and multiplication combine supports already present.  Composition
changes the scale itself: replacing $X$ by $G(X)$ transforms both
$X^{-1}$ and $\log X$, and the transformed generators may introduce
fractional powers, new logarithmic levels, or entirely new exponential
monomials.  In the full transseries field, composition is defined for a large
positive inner transseries and is constrained by summability; Edgar's
composition paper gives a detailed treatment \cite{EdgarComposition}.  In the
polynomial--logarithmic setting, the most important and cleanest case is a
near-identity argument.

\section{Near-identity arguments}
Let
\begin{equation}\label{eq:G-near-id}
 G(X)=X+\Delta(X)=X(1+U(X)),
 \qquad U=t\Delta.
\end{equation}
Assume
\begin{equation}\label{eq:U-small}
 U\in t\mathcal C_L[[t]],
\end{equation}
so $U\precdom1$.  The two formal generators transform exactly as
\begin{equation}\label{eq:generator-substitution}
 t=\frac1X\longmapsto
 \frac1G=t(1+U)^{-1},
 \qquad
 L=\log X\longmapsto
 \log G=L+\log(1+U).
\end{equation}
Both right-hand sides belong to $\PLfield$ and are computable by geometric
and logarithmic series.

\begin{definition}[Near-identity composition]
For
\[
 F(X)=\sum_{n\ge n_0}p_n(L)t^n
\]
and $G$ satisfying \eqref{eq:G-near-id}--\eqref{eq:U-small}, define
\begin{equation}\label{eq:composition-generator-formula}
 (F\compose G)(X)
 =\sum_{n\ge n_0}
 t^n(1+U)^{-n}
 p_n\!\left(L+\log(1+U)\right).
\end{equation}
\end{definition}

\begin{theorem}[Well-definedness at finite order]\label{thm:composition-well-defined}
For every outer cutoff $N$, the truncation of \eqref{eq:composition-generator-formula}
through $t^N$ depends on only finitely many input blocks and finitely many
terms of each binomial, logarithmic, and polynomial Taylor expansion.
Therefore near-identity composition is a well-defined operation on
$\PLfield$.
\end{theorem}
\begin{proof}
Because $U=\FormalBigOAt{t}{t}$, the coefficient of $t^r$ in $(1+U)^{-n}$ or
$\log(1+U)$ depends on finitely many powers of $U$.  A coefficient polynomial
$p_n$ has finite degree, so its substitution at
$L+\log(1+U)$ is a finite Taylor polynomial in the logarithmic increment.
Finally, only blocks $n\le N$ can affect the output through $t^N$.
\end{proof}

\section{A direct calculation}
Take
\[
 F(X)=\log X+\frac1X=L+t,
 \qquad
 G(X)=X+\log X=X+L.
\]
Here $U=Lt$.  The transformed generators are
\begin{align*}
 \log G
 &=L+\log(1+Lt)\\
 &=L+Lt-\frac12L^2t^2+\frac13L^3t^3
   -\frac14L^4t^4+\FormalBigOAt{t^5}{t},\\
 \frac1G
 &=t(1+Lt)^{-1}\\
 &=t-Lt^2+L^2t^3-L^3t^4+\FormalBigOAt{t^5}{t}.
\end{align*}
Therefore
\begin{equation}\label{eq:composition-example}
 F\compose G
 =L+(L+1)t-\left(L+\frac12L^2\right)t^2
 +\left(L^2+\frac13L^3\right)t^3
 -\left(L^3+\frac14L^4\right)t^4
 +\FormalBigOAt{t^5}{t}.
\end{equation}
One may check the same result by Taylor composition in the next chapter.

\section{Composition as an algebra homomorphism}
For fixed admissible $G$, substitution respects the ring operations:
\begin{align}\label{eq:composition-homomorphism}
 (F_1+F_2)\compose G&=(F_1\compose G)+(F_2\compose G),\\
 (F_1F_2)\compose G&=(F_1\compose G)(F_2\compose G).
\end{align}
It also respects reciprocal whenever the factor is nonzero:
\[
 (F^{-1})\compose G=(F\compose G)^{-1}.
\]
These statements follow because \eqref{eq:generator-substitution} defines a
homomorphism from the algebra generated by $t,L$ and summable expansions into
itself.

\section{Associativity}
Suppose $G$ and $H$ are near the identity and all indicated substitutions
remain in the same admissible class.  Then
\begin{equation}\label{eq:composition-associativity}
 (F\compose G)\compose H=F\compose(G\compose H).
\end{equation}
Indeed, both sides replace $X$ by $G(H(X))$.  On the generators,
\[
 \frac1X\mapsto\frac1{G(H(X))},
 \qquad
 \log X\mapsto\log(G(H(X))),
\]
and homomorphic extension gives equality.  The qualification about
admissibility matters in larger Hahn fields: a syntactically meaningful
nested expression need not define a summable support without hypotheses.

\section{Practical generator-substitution algorithm}
To compute $F\compose G$ through block $N$:
\begin{enumerate}
\item Form $U=t(G-X)$ through block $N+1$.
\item Compute $R=(1+U)^{-1}$ and $V=\log(1+U)$ through the required order.
\item For every input block $t^np_n(L)$ with $n\le N$, compute
      $t^nR^np_n(L+V)$.
\item Sum and truncate after each multiplication.
\end{enumerate}
If negative $n$ occur, $R^n$ is an ordinary positive power for $-n>0$.
Guard terms are needed because an initial $t^{-1}$ block can shift a term of
$R^n$ one place toward the retained range.

\begin{algorithm}[H]
\caption{Near-identity composition through outer block $N$}
\label{alg:generator-composition}
\KwIn{$F=\sum_{n=n_0}^{N}p_nt^n$, $G=X+\Delta$, target $N$}
\KwOut{$\trunc{F\compose G}{N}$}
$U\leftarrow t\Delta$\;
$R\leftarrow\trunc{(1+U)^{-1}}{N-n_0}$\;
$V\leftarrow\trunc{\log(1+U)}{N-n_0}$\;
$H\leftarrow0$\;
\For{$n\leftarrow n_0$ \KwTo $N$}{
  $H\leftarrow H+t^nR^np_n(L+V)$\;
  $H\leftarrow\trunc{H}{N}$\;
}
\Return{$H$}\;
\end{algorithm}

\begin{pitfallbox}
Never substitute $L\mapsto\log G$ while leaving $t$ unchanged.  The symbols
$t$ and $L$ are linked by $t=X^{-1}$ and $L=\log X$; composition must transform
both.  Updating only one generator is a common source of plausible but false
coefficients.
\end{pitfallbox}

\begin{summarybox}
For $G=X(1+U)$ with $U\precdom1$, composition is exact generator
substitution:
$t\mapsto t(1+U)^{-1}$ and
$L\mapsto L+\log(1+U)$.  This construction is coefficientwise finite,
respects algebraic operations, and makes associativity transparent.
\end{summarybox}


\chapter{Taylor composition and low-order coefficient laws}
\label{ch:composition-taylor}

\section{The formal Taylor formula}
Let $G=X+\Delta$ with $\Delta\in\mathcal C_L[[t]]$, so
$t\Delta\precdom1$.  Since differentiation raises outer order, the formal
Taylor family
\begin{equation}\label{eq:taylor-composition}
 F(X+\Delta)
 =\sum_{r\ge0}\frac{\Delta^r}{r!}\D^rF(X)
\end{equation}
is coefficientwise summable.

\begin{theorem}[Taylor composition in the one-logarithm field]
\label{thm:taylor-composition}
For $F\in\PLfield$ and $G=X+\Delta$ as above,
\begin{equation}\label{eq:taylor-composition-theorem}
 F\compose G
 =\sum_{r\ge0}\frac{\Delta^r}{r!}\D^rF.
\end{equation}
At any fixed outer order $N$, only finitely many $r$ contribute.
\end{theorem}
\begin{proof}
For a block $P(L)t^n$, \cref{thm:repeated-derivative} shows that
$\D^r(Pt^n)$ has outer order $n+r$.  Since $\nu_t\Delta\ge0$, the $r$th
Taylor term has order at least $n+r$, so only $r\le N-n$ can affect block
$N$.  It remains to identify the sum with generator substitution.  Applied
to $t=X^{-1}$, the Taylor series is
\[
 \sum_{r\ge0}\frac{\Delta^r}{r!}\D^rt
 =t\sum_{r\ge0}(-t\Delta)^r
 =\frac{t}{1+t\Delta}=\frac1G.
\]
Applied to $L=\log X$, it is
\[
 L+\sum_{r\ge1}\frac{(-1)^{r-1}}r(t\Delta)^r
 =L+\log(1+t\Delta)=\log G.
\]
The equality then extends to polynomials, products, inverses, and summable
series.
\end{proof}

\section{When Taylor's formula is not justified}
The condition is not merely that $\Delta$ tends to zero numerically.  What is
needed is that successive differentiated terms decrease in the transseries
ordering.  In a larger exponential field, take $F(X)=\exp(\exp X)$ and
$\Delta=\exp(-X/2)$.  Although $\Delta\to0$,
\[
 \Delta F'(X)=\exp(-X/2)\exp X\exp(\exp X)
 =\exp(X/2)F(X),
\]
which is larger than $F(X)$.  Thus the naive Taylor hierarchy fails.  The
polynomial--logarithmic near-identity class avoids this problem because each
derivative supplies a factor $X^{-1}$ that dominates polynomial logarithmic
growth.

\section{Low-order composition formulas}
Write
\begin{align}\label{eq:F-G-coeff-notation}
 F(X)&=X+f_0(L)+t f_1(L)+t^2f_2(L)+t^3f_3(L)+\FormalBigOAt{t^4}{t},\\
 G(X)&=X+g_0(L)+t g_1(L)+t^2g_2(L)+t^3g_3(L)+\FormalBigOAt{t^4}{t}.
\end{align}
Let
\[
 H=F\compose G
 =X+h_0+t h_1+t^2h_2+t^3h_3+\FormalBigOAt{t^4}{t}.
\]
Primes below mean derivatives with respect to $L$.

\begin{proposition}[Composition through three inverse-power blocks]
\label{prop:composition-first-blocks}
The coefficients are
\begin{align}
 h_0={}&f_0+g_0,\label{eq:h0}\\
 h_1={}&g_1+f_1+g_0f_0',\label{eq:h1}\\
 h_2={}&g_2+f_2+g_1f_0'
 +g_0(f_1'-f_1)
 +\frac{g_0^2}{2}(f_0''-f_0'),\label{eq:h2}\\
 h_3={}&g_3+f_3+g_2f_0'
 +g_1(f_1'-f_1)+g_0(f_2'-2f_2)\notag\\
 &+g_0g_1(f_0''-f_0')
 +\frac{g_0^2}{2}(f_1''-3f_1'+2f_1)\notag\\
 &+\frac{g_0^3}{6}(f_0'''-3f_0''+2f_0').
 \label{eq:h3}
\end{align}
\end{proposition}
\begin{proof}
Set $U=t(G-X)$.  Through $t^3$,
\begin{align*}
 U&=g_0t+g_1t^2+g_2t^3+\cdots,\\
 \log(1+U)
 &=g_0t+\left(g_1-\frac12g_0^2\right)t^2
 +\left(g_2-g_0g_1+\frac13g_0^3\right)t^3+\cdots,\\
 (1+U)^{-1}
 &=1-g_0t+(g_0^2-g_1)t^2+\cdots,\\
 (1+U)^{-2}
 &=1-2g_0t+\cdots.
\end{align*}
Insert these in \eqref{eq:composition-generator-formula}, Taylor-expand the
$f_j$ in the logarithmic increment, and collect blocks.  The displayed
formulas result.
\end{proof}

\section{A check using a simple shift}
If $G=X+c$ with constant $c$, then $g_0=c$ and all later $g_j$ vanish.  The
formulas reduce to
\begin{align*}
 h_0&=f_0+c,\\
 h_1&=f_1+cf_0',\\
 h_2&=f_2+c(f_1'-f_1)+\frac{c^2}{2}(f_0''-f_0'),
\end{align*}
which agrees with the ordinary Taylor expansion in $c$ because
\[
 \D f_0(L)=tf_0',
 \qquad
 \D^2 f_0(L)=t^2(f_0''-f_0').
\]

\section{Noncommutativity}
Composition is associative but not commutative.  Let
\[
 F=X+L,
 \qquad
 G=X+L^2.
\]
Then
\begin{align*}
 F\compose G
 &=X+L^2+L+L^2t-\frac12L^4t^2+\cdots,\\
 G\compose F
 &=X+L+L^2+2L^2t+(L^2-L^3)t^2+\cdots.
\end{align*}
Therefore
\begin{equation}\label{eq:composition-noncommutative}
 F\compose G-G\compose F=-L^2t+\FormalBigOAt{t^2}{t}.
\end{equation}
At first order, the commutator is governed by the vector-field bracket
$f_0g_0'-g_0f_0'$.  This is the beginning of the Lie-theoretic structure
behind formal iteration.

\section{Taylor form versus generator form}
The generator formula is often best when $F$ is explicitly stored as blocks
$p_n(L)t^n$.  The Taylor formula is best when derivatives of $F$ are already
available or when one seeks universal coefficient identities.  They have the
same formal content, but different computational profiles:
\begin{center}
\begin{tabularx}{0.94\textwidth}{@{}lXX@{}}
\toprule
 & Generator substitution & Taylor composition \\
\midrule
Core data & $(1+U)^{-1}$ and $\log(1+U)$ & $\D^rF$ and $\Delta^r/r!$ \\
Best for & sparse coefficient blocks & symbolic universal formulas \\
Truncation & powers of $U$ & derivative order \\
Main risk & forgetting one generator & using an invalid smallness condition \\
\bottomrule
\end{tabularx}
\end{center}

\begin{summarybox}
Near the identity, formal Taylor composition is a theorem because every
derivative raises the inverse-power order.  It is equivalent to exact
substitution of $1/X$ and $\log X$.  The first coefficient laws are triangular
and expose both noncommutativity and the mechanism needed for series reversal.
\end{summarybox}


\chapter{The near-identity composition group}
\label{ch:composition-group}

\section{The group of tangent-to-identity transseries at infinity}
Define
\begin{equation}\label{eq:near-id-group}
 \mathcal G_{\mathrm{itLaur}}
 =\SetOf{F=X+A(X):A\in\mathcal C_L[[t]]}.
\end{equation}
Every such $F$ satisfies $F/X=1+tA=1+\LittleO(1)$ formally.  The identity is $X$.
Closure under composition follows from the preceding chapters.

\begin{theorem}[Near-identity group]\label{thm:near-id-group}
The set $\mathcal G_{\mathrm{itLaur}}$ is a group under composition.  Every
$F=X+A$ has a unique inverse of the form $X+B$ with
$B\in\mathcal C_L[[t]]$.
\end{theorem}
The proof will be constructive in Part~\ref{part:reversion}: fixed-point
iteration gains one outer block at each step.  The result is a concrete
one-logarithm instance of the compositional group of large positive
transseries treated by Edgar \cite{EdgarComposition}.

\section{Filtration by contact order}
For $F=X+A$ and $G=X+B$, say that they agree through contact order $r$ when
$A-B=\FormalBigOAt{t^r}{t}$.  If $H=X+C$ is fixed, then
\begin{equation}\label{eq:composition-continuity-left}
 F-G=\FormalBigOAt{t^r}{t}
 \quad\Longrightarrow\quad
 H\compose F-H\compose G=\FormalBigOAt{t^r}{t},
\end{equation}
and, because $H'=1+\FormalBigOAt{t}{t}$,
\begin{equation}\label{eq:composition-continuity-right}
 F\compose H-G\compose H=\FormalBigOAt{t^r}{t}.
\end{equation}
This filtration makes composition continuous and allows infinite coefficient
constructions by stabilization.

\section{The first commutator block}
Let $F=X+f(L)$ and $G=X+g(L)$ with no inverse-power corrections.  Using
\eqref{eq:h1},
\begin{align*}
 F\compose G&=X+f+g+t\,g f'+\FormalBigOAt{t^2}{t},\\
 G\compose F&=X+f+g+t\,f g'+\FormalBigOAt{t^2}{t}.
\end{align*}
Hence
\begin{equation}\label{eq:first-commutator}
 F\compose G-G\compose F
 =t(gf'-fg')+\FormalBigOAt{t^2}{t}.
\end{equation}
The bilinear operation $[f,g]=gf'-fg'$ is the usual bracket of one-dimensional
vector fields after the change of independent variable $L=\log X$.

\section{Iterates and fractional iteration}
Integer iterates $F^{[n]}$ are defined by repeated composition.  Because
$\mathcal G_{\mathrm{itLaur}}$ is filtered and pro-unipotent at finite truncation, one can
formally introduce the substitution operator
\[
 \mathcal S_F:H\longmapsto H\compose F
\]
and, when $\mathcal S_F-I$ raises the filtration, define
\begin{equation}\label{eq:substitution-log}
 \log\mathcal S_F
 =\sum_{r\ge1}\frac{(-1)^{r+1}}r(\mathcal S_F-I)^r.
\end{equation}
Then a formal fractional iterate may be written
\[
 \mathcal S_F^s=\exp\bigl(s\log\mathcal S_F\bigr),
 \qquad
 F^{[s]}=X\,\mathcal S_F^s.
\]
Developing this construction fully leads to iterative logarithms and Abel
coordinates, beyond the central scope here; Edgar's work on fractional
iteration provides a transseries treatment \cite{EdgarFractional}.  The key
point for the present article is that the same filtration that legitimizes
composition also legitimizes recursive and operator-logarithmic constructions.

\section{Coefficient stabilization under iteration}
Suppose a map $\Phi$ satisfies
\begin{equation}\label{eq:contraction-general}
 U-V=\FormalBigOAt{t^r}{t}
 \quad\Longrightarrow\quad
 \Phi(U)-\Phi(V)=\FormalBigOAt{t^{r+1}}{t}.
\end{equation}
Then the iterates $U_{n+1}=\Phi(U_n)$ stabilize one additional block per step:
$U_{n+1}-U_n=\FormalBigOAt{t^n}{t}$ after a suitable initialization.  This is the
formal analogue of a contraction mapping, but no numerical norm is needed.
The ordered filtration supplies the notion of increasing precision.

\begin{summarybox}
Near-identity polynomial--logarithmic transseries form a noncommutative group.
Their $t$-adic contact filtration makes composition continuous, identifies the
first commutator with a vector-field bracket, and supports fixed points,
inverses, and even formal fractional iteration.
\end{summarybox}


\chapter{Scale-changing composition and deeper logarithms}
\label{ch:composition-scale}

\section{A general leading inner scale}
Suppose
\begin{equation}\label{eq:G-general-leading}
 G(X)=cX^\rho L^\sigma(1+S),
 \qquad c\ne0,
 \qquad \rho>0,
 \qquad S\precdom1.
\end{equation}
On a real positive ray take $c>0$; in the complex plane choose branches.
Then
\begin{align}
 \frac1G&=c^{-1}t^\rho L^{-\sigma}(1+S)^{-1},
 \label{eq:scale-t-transform}\\
 \log G&=\rho L+\sigma\log L+\log c+\log(1+S).
 \label{eq:scale-L-transform}
\end{align}
These two formulas completely determine $F\compose G$.

\section{What new scales are forced?}
Equation \eqref{eq:scale-t-transform} may require fractional powers $t^\rho$.
Equation \eqref{eq:scale-L-transform} introduces $\log L$ if $\sigma\ne0$.
If the coefficients of $F$ already involve $\log L$, composition may create
$\log(\rho L+\sigma\log L+\cdots)$ and hence another nested expansion.  The
correct ambient field is chosen by examining these transformed generators,
not by guessing from the appearance of $F$ alone.

\begin{example}[A scale-changing substitution]
Let
\[
 F(X)=X+\log X+X^{-1},
 \qquad
 G(X)=3X^2L^{-1}.
\]
Then
\begin{align*}
 F\compose G
 &=3X^2L^{-1}
 +\log(3X^2L^{-1})
 +\frac{L}{3X^2}\\
 &=3X^2L^{-1}+2L-\log L+\log3+\frac13Lt^2.
\end{align*}
A second logarithmic level appears exactly, not merely as a remainder.
\end{example}

\section{Composition with a logarithm}
If $G=\log X=L$, then
\[
 t\mapsto L^{-1},
 \qquad
 L\mapsto\log L.
\]
Thus
\[
 \left(\sum_np_n(L)t^n\right)\compose\log
 =\sum_np_n(\log L)L^{-n}.
\]
This lowers the exponential height but increases logarithmic depth.  Repeated
composition with $\log$ produces the hierarchy
$L_1,L_2,L_3,\ldots$.

\section{Composition with a power}
For $G=X^\rho$ with $\rho>0$,
\[
 t\mapsto t^\rho,
 \qquad
 L\mapsto\rho L.
\]
Therefore
\begin{equation}\label{eq:composition-power}
 F(X^\rho)=\sum_n t^{\rho n}p_n(\rho L).
\end{equation}
This operation stays within an integer lattice only when $\rho$ maps the
chosen exponent group into itself.  A rational $\rho=p/q$ is handled by the
generator $X^{-1/q}$.

\section{Failure when the inner argument is not large}
The formal logarithm $\log G$ requires a leading monomial and a normalized
small remainder.  If $G$ oscillates across zero or has incomparable dominant
terms in an analytic problem, a single branch may not exist on the domain of
interest.  Likewise, if $G$ is small rather than large, composing an expansion
valid at infinity with $G$ has no asymptotic meaning.  Formal syntax cannot
replace domain and branch analysis.

\section{An admissibility checklist}
Before computing $F\compose G$, determine:
\begin{enumerate}
\item Is $G$ in the asymptotic regime for which $F$ is expanded?
\item What is the leading monomial of $G$?
\item Can $G$ be factored as leading monomial times $1+S$ with $S\precdom1$?
\item What do $1/G$ and $\log G$ look like?
\item Do they remain in the chosen exponent lattice and logarithmic depth?
\item In the complex plane, which branches of powers and logarithms are used?
\item Is the resulting family of supports summable at the desired cutoff?
\end{enumerate}

\begin{summarybox}
Scale-changing composition is controlled by the images of the generators.
For $G\sim cX^\rho L^\sigma$, one has
$1/G\sim c^{-1}X^{-\rho}L^{-\sigma}$ and
$\log G=\rho L+\sigma\log L+\log c+\LittleO(1)$.  Fractional powers and deeper
logarithms are therefore often mathematically unavoidable.
\end{summarybox}

\part{Series reversal and asymptotic inversion}
\label{part:reversion}

\chapter{Fixed-point construction of the compositional inverse}
\label{ch:reversion-fixed}

\section{The inverse equation}
Let
\begin{equation}\label{eq:F-near-id-reversion}
 F(X)=X+A(X),
 \qquad A\in\mathcal C_L[[t]].
\end{equation}
We seek $G=F\compinv$ in the form
\[
 G(Y)=Y+B(Y).
\]
The equation $F(G(Y))=Y$ is equivalent to
\begin{equation}\label{eq:inverse-fixed-B}
 B(Y)=-A(Y+B(Y)).
\end{equation}
Thus $B$ is a fixed point of
\begin{equation}\label{eq:Phi-inverse}
 \Phi(B)=-A\compose(Y+B).
\end{equation}
The variable is written $Y$ during inversion to separate the input of the
inverse from the original variable $X$; after the calculation one may rename
it $X$.

\section{The one-block contraction}
The derivative of $A$ has outer order at least one.  Therefore changing its
argument by a block-$r$ error changes the value only at block $r+1$ or later.

\begin{lemma}[Filtered Lipschitz gain]\label{lem:inverse-contraction}
If $B-C=\FormalBigOAt{t^r}{t}$, with $B,C\in\mathcal C_L[[t]]$, then
\begin{equation}\label{eq:inverse-contraction}
 A(Y+B)-A(Y+C)=\FormalBigOAt{t^{r+1}}{t}.
\end{equation}
\end{lemma}
\begin{proof}
Taylor-expand around $Y+C$:
\[
 A(Y+B)-A(Y+C)
 =\sum_{j\ge1}\frac{(B-C)^j}{j!}
 \bigl(\D^jA\bigr)(Y+C).
\]
The derivative $\D^jA$ has outer order at least $j$.  The $j$th term therefore
has order at least $jr+j\ge r+1$ for $r\ge0$.
\end{proof}

\begin{theorem}[Existence and uniqueness by coefficient stabilization]
\label{thm:fixed-inverse}
For every $F$ of the form \eqref{eq:F-near-id-reversion}, the iteration
\begin{equation}\label{eq:fixed-inverse-iteration}
 B_0=0,
 \qquad
 B_{n+1}=-A(Y+B_n)
\end{equation}
stabilizes coefficientwise.  Its limit $B$ is the unique solution of
\eqref{eq:inverse-fixed-B}.  Consequently $F$ has a unique near-identity
compositional inverse $G=Y+B$.
\end{theorem}
\begin{proof}
The first difference $B_1-B_0=-A(Y)$ has outer order at least zero.
By \cref{lem:inverse-contraction},
\[
 B_{n+1}-B_n=\FormalBigOAt{t^n}{t}.
\]
Hence each fixed block eventually stops changing, which defines a formal
limit $B$.  Composition is continuous in the filtration, so the limit
satisfies the fixed-point equation.  If $B$ and $C$ were two solutions and
their first difference occurred at block $r$, then
$B-C=\Phi(B)-\Phi(C)=\FormalBigOAt{t^{r+1}}{t}$, a contradiction.
\end{proof}

\section{Why a bad rearrangement may fail}
For Lambert inversion one can rearrange the equation in more than one way.
If $X=\LambertW\EulerE^\LambertW$, then $\LambertW=\log X-\log \LambertW$, which is contractive after the dominant
term is removed.  The equivalent equation $Q=\EulerE^{-Q}-\log X$ for a correction
$Q$ is not contractive at the natural starting point: its exponential can
create a term as large as $X$.  Edgar uses this example to emphasize that a
formal fixed-point scheme must increase asymptotic precision, not merely be
algebraically equivalent \cite{EdgarBeginners}.

\section{Stopping after a requested order}
To compute through $t^N$, it is enough to perform $N+1$ fixed-point steps from
$B_0=0$, although some blocks may stabilize sooner.  After every step one may
truncate at $t^{N+1}$ because the next application of $A$ cannot bring a
discarded block back to an earlier order.

\begin{algorithm}[H]
\caption{Fixed-point compositional inversion}
\label{alg:fixed-inverse}
\KwIn{$F=X+A$, target outer block $N$}
\KwOut{$G=\trunc{F\compinv}{N}$}
$B\leftarrow0$\;
\For{$j\leftarrow0$ \KwTo $N$}{
  $B\leftarrow\trunc{-\,A\compose(Y+B)}{N}$\;
}
\Return{$Y+B$}\;
\end{algorithm}

\section{Left and right inverse}
The fixed-point construction gives $F\compose G=Y$.  In a group this already
implies $G\compose F=X$, but it is useful to see why directly.  Construct a
right inverse $H$ of $G$ by the same theorem.  Then
\[
 F=F\compose(G\compose H)=(F\compose G)\compose H=H.
\]
Thus $G\compose F=G\compose H=Y$.  The inverse is two-sided.

\begin{summarybox}
For $F=X+A$ with $A=\LittleO(X)$ in the polynomial--logarithmic scale, the equation
$B=-A(Y+B)$ is contractive in outer order because $A'=\FormalBigOAt{t}{t}$.  Simple
iteration stabilizes one new block at every step and proves both existence and
uniqueness of the compositional inverse.
\end{summarybox}


\chapter{Triangular coefficient reversal}
\label{ch:reversion-triangular}

\section{Solving one block at a time}
Use the notation
\begin{align*}
 F(X)&=X+f_0+t f_1+t^2f_2+t^3f_3+\cdots,\\
 G(Y)&=Y+g_0+t g_1+t^2g_2+t^3g_3+\cdots,
\end{align*}
where all coefficients are functions of $L=\log Y$.  Set $F\compose G=Y$ in
\cref{prop:composition-first-blocks}.  Each equation is linear in the new
unknown $g_n$, with coefficient one.  This is the triangular structure of
series reversal.

The first three inverse blocks are
\begin{align}
 g_0={}&-f_0,\label{eq:inverse-g0}\\
 g_1={}&-f_1-g_0f_0'= -f_1+f_0f_0',\label{eq:inverse-g1}\\
 g_2={}&-f_2-g_1f_0'-g_0(f_1'-f_1)
 -\frac{g_0^2}{2}(f_0''-f_0').
 \label{eq:inverse-g2}
\end{align}
The formula for $g_3$ follows by negating every term of \eqref{eq:h3} except
$g_3$.  At higher orders one normally lets a composition routine generate the
known residual and then sets the next inverse block equal to its negative.

\section{Universal triangular algorithm}
Suppose $G_{<n}=Y+\sum_{j=0}^{n-1}g_jt^j$ already satisfies
\[
 F\compose G_{<n}=Y+r_nt^n+\FormalBigOAt{t^{n+1}}{t}.
\]
Adding $g_nt^n$ to $G$ changes $F(G)$ by
\[
 F'(G_{<n})g_nt^n+\FormalBigOAt{t^{n+1}}{t}
 =g_nt^n+\FormalBigOAt{t^{n+1}}{t},
\]
because $F'=1+\FormalBigOAt{t}{t}$.  Hence
\begin{equation}\label{eq:triangular-next-block}
 g_n=-r_n.
\end{equation}
This gives an especially simple implementation: compose the current inverse,
read the first nonzero residual block, and cancel it.

\begin{algorithm}[H]
\caption{Residual-cancellation series reversal}
\label{alg:triangular-reversal}
\KwIn{$F=X+A$, target outer block $N$}
\KwOut{$G=\trunc{F\compinv}{N}$}
$G\leftarrow Y$\;
\For{$n\leftarrow0$ \KwTo $N$}{
  $R\leftarrow\trunc{F\compose G-Y}{n}$\;
  $r_n\leftarrow[t^n]R$\;
  $G\leftarrow G-r_nt^n$\;
}
\Return{$G$}\;
\end{algorithm}

\section{A seven-parameter example}
Consider
\begin{equation}\label{eq:parametric-F}
 F(X)=X+aL+b+(cL+d)t+(eL^2+fL+g)t^2+\FormalBigOAt{t^3}{t}.
\end{equation}
The inverse has the form
\[
 G(Y)=Y+g_0(L)+g_1(L)t+g_2(L)t^2+\FormalBigOAt{t^3}{t},
\]
with
\begin{align}
 g_0={}&-(aL+b),\label{eq:param-g0}\\
 g_1={}&(a^2-c)L+(ab-d),\label{eq:param-g1}\\
 g_2={}&\left(\frac{a^3}{2}-ac-e\right)L^2\notag\\
 &+\left(-a^3+a^2b+2ac-ad-bc-f\right)L\notag\\
 &-a^2b+\frac{ab^2}{2}+ad+bc-bd-g.
 \label{eq:param-g2}
\end{align}
These formulas are useful as unit tests for symbolic implementations: random
numerical values of the parameters can be substituted and the composition
checked through $t^2$.

\section{A one-parameter perturbation of Lambert inversion}
Let
\begin{equation}\label{eq:F-log-q}
 F(X)=X+\log X+\frac{q}{X}=X+L+qt.
\end{equation}
Substituting $a=1$, $b=0$, $c=0$, $d=q$, and $e=f=g=0$ into the preceding
formulas gives
\begin{equation}\label{eq:inverse-log-q}
 F\compinv(Y)
 =Y-L+(L-q)t
 +\left(\frac12L^2-(q+1)L+q\right)t^2
 +\FormalBigOAt{t^3}{t}.
\end{equation}
At $q=0$ this begins the Wright omega, hence Lambert-$\LambertW$, expansion developed
in Part~\ref{part:lambert}.

\section{Degree growth}
If $\deg f_0=d$, repeated composition can make $\deg g_n$ grow roughly
linearly with $n$.  In the Lambert case $f_0=L$, one has exactly
$\deg g_n=n$ for $n\ge1$.  The triangular algorithm never requires an a
priori degree guess: coefficient polynomials emerge from arithmetic.  For a
preallocated dense representation, however, the degree bound supplied by
composition formulas can substantially reduce memory.

\begin{summarybox}
Series reversal is triangular because $F'=1+\FormalBigOAt{t}{t}$.  Once an approximate
inverse is correct through block $n-1$, the block-$n$ residual is canceled by
subtracting exactly that coefficient from the inverse.  Universal low-order
formulas and residual cancellation provide complementary symbolic methods.
\end{summarybox}


\chapter{Newton iteration for compositional inversion}
\label{ch:reversion-newton}

\section{The Newton map}
For a current approximation $G(Y)$ to $F\compinv(Y)$, define the residual
\[
 R=F\compose G-Y.
\]
Newton's correction is
\begin{equation}\label{eq:newton-compositional}
 G_{\mathrm{new}}
 =G-\frac{F\compose G-Y}{F'\compose G}.
\end{equation}
The denominator is a unit because $F'=1+\FormalBigOAt{t}{t}$.

\begin{theorem}[Precision doubling]\label{thm:newton-doubling}
If $R=\FormalBigOAt{t^m}{t}$, then the residual after
\eqref{eq:newton-compositional} is at least
$\FormalBigOAt{t^{2m+2}}{t}$ when $m\ge0$.  In particular, the number of correct outer
blocks approximately doubles at each step.
\end{theorem}
\begin{proof}
Put $C=-R/(F'\compose G)$.  Then $C=\FormalBigOAt{t^m}{t}$.  Taylor expansion gives
\begin{align*}
 F\compose(G+C)-Y
 &=R+(F'\compose G)C
 +\sum_{j\ge2}\frac{C^j}{j!}(F^{(j)}\compose G).
\end{align*}
The first two terms cancel.  Since $F^{(j)}=A^{(j)}$ has outer order at least
$j$, the $j$th remaining term has order at least $jm+j$.  The minimum occurs
at $j=2$ and is $2m+2$.
\end{proof}

\section{Newton applied to \texorpdfstring{$X+\log X$}{X+log X}}
Let $F(X)=X+\log X$ and start with $G_0=Y$.  Then
\[
 F(G_0)-Y=L,
 \qquad
 F'(G_0)=1+t.
\]
The first Newton iterate is
\begin{align*}
 G_1
 &=Y-\frac{L}{1+t}\\
 &=Y-L+Lt-Lt^2+Lt^3-\cdots.
\end{align*}
It is already correct through block $t^1$.  A second Newton step produces the
correct expansion through at least $t^5$ if all intermediate arithmetic is
carried with sufficient guard precision.  Fixed-point iteration would require
one step per newly stabilized block.

\section{Truncated Newton algorithm}
\begin{algorithm}[H]
\caption{Newton compositional inversion}
\label{alg:newton-composition}
\KwIn{$F=X+A$, target outer block $N$}
\KwOut{$G=\trunc{F\compinv}{N}$}
$G\leftarrow Y$; $p\leftarrow0$\;
\While{$p\le N$}{
  choose $q\leftarrow\min(N+1,2p+2)$\;
  $R\leftarrow\trunc{F\compose G-Y}{q-1}$\;
  $J\leftarrow\trunc{F'\compose G}{q-1}$\;
  $G\leftarrow\trunc{G-R/J}{q-1}$\;
  $p\leftarrow q$\;
}
\Return{$G$}\;
\end{algorithm}
The variable $p$ records a conservative residual order.  In practice one
checks the residual after every update rather than relying solely on a
predicted precision.

\section{Cost comparison}
Let $M(N)$ denote the cost of truncated multiplication to order $N$, and let
$C(N)$ denote the cost of composition.  A triangular method performs roughly
$N$ compositions at growing precision, while Newton performs $\BigO(\log N)$
compositions and divisions at geometrically increasing precision.  With fast
arithmetic this makes Newton asymptotically superior.  At low order, however,
the triangular recurrence is simpler, exposes exact coefficient structure,
and is easier to debug.

\section{Hybrid strategy}
A practical symbolic system often uses:
\begin{enumerate}
\item dominant balance to obtain the leading term;
\item two or three triangular blocks to establish the correct scale and
      branches;
\item Newton iteration to reach high order;
\item a final residual computation as certification.
\end{enumerate}
The methods are not competitors but stages of one robust workflow.

\begin{pitfallbox}
Newton's formula is only as accurate as its denominator and composition.
Computing $F'\compose G$ to the final cutoff but computing the residual at
lower precision can silently lose blocks after division.  Every Newton step
must carry a common precision context with guard terms determined by leading
orders.
\end{pitfallbox}

\begin{summarybox}
Compositional Newton iteration subtracts the residual divided by the composed
derivative.  Because higher derivatives of the correction $A$ begin at
progressively later inverse-power blocks, Newton turns an order-$m$ residual
into one of order at least $2m+2$.  It is the method of choice for high-order
computation.
\end{summarybox}


\chapter{The transserial Lagrange--B\"urmann formula}
\label{ch:reversion-lagrange}

\section{A closed all-orders formula}
Fixed-point and Newton iteration are recursive.  Lagrange--B\"urmann inversion
packages every correction into one explicit family.

\begin{theorem}[Lagrange--B\"urmann reversion at infinity]
\label{thm:lagrange-transserial}
Let $F(X)=X+A(X)$ with $A\in\mathcal C_L[[t]]$.  Then
\begin{equation}\label{eq:lagrange-main}
 F\compinv(Y)
 =Y+\sum_{n\ge1}\frac{(-1)^n}{n!}
 \D_Y^{\,n-1}\bigl(A(Y)^n\bigr).
\end{equation}
The family is coefficientwise summable in the polynomial--logarithmic
filtration: the block $t^N$ receives contributions only from $n\le N+1$.
\end{theorem}

\begin{proof}
Introduce a formal parameter $z$ and solve
\[
 U=z\phi(Y+U),
 \qquad \phi=-A.
\]
The ordinary formal Lagrange--B\"urmann theorem gives
\begin{equation}\label{eq:lagrange-parameter}
 U=\sum_{n\ge1}\frac{z^n}{n!}
 \D_Y^{\,n-1}\bigl(\phi(Y)^n\bigr).
\end{equation}
For completeness, its coefficient form is
\[
 [z^n]H(Y+U)
 =\frac1n[w^{n-1}]H'(Y+w)\phi(Y+w)^n.
\]
This follows by the formal residue substitution
$z=w/\phi(Y+w)$; an integration by parts reduces the resulting residue to the
displayed expression.  Taking $H(u)=u$ gives
\eqref{eq:lagrange-parameter}.  Now set $z=1$.  Since $A^n$ has outer order at
least zero and $n-1$ derivatives raise outer order by at least $n-1$, only
$n\le N+1$ can affect block $N$.  Thus evaluation at $z=1$ is legitimate
coefficientwise and yields \eqref{eq:lagrange-main}.
\end{proof}

\section{The first universal terms}
Expanding \eqref{eq:lagrange-main} gives
\begin{align}
 F\compinv(Y)={}&Y-A+AA'
 -A(A')^2-\frac12A^2A''\notag\\
 &+A(A')^3+\frac32A^2A'A''+\frac16A^3A'''\notag\\
 &-A(A')^4-3A^2(A')^2A''-\frac12A^3(A'')^2
 -\frac23A^3A'A'''-\frac1{24}A^4A^{(4)}+\cdots.
 \label{eq:lagrange-first-terms}
\end{align}
All functions and derivatives on the right are evaluated at $Y$.  The fourth
line follows from $-(1/120)\D^4(A^5)$; its partition coefficients are a
specialization of the Bell-polynomial formulas reviewed in
Appendix~\ref{app:bell}.

\section{Lagrange formula for an observable}
The more general form is often as useful as the inverse itself.

\begin{theorem}[Composed observable]
For any differentiable formal transseries $H$ for which the family is
summable,
\begin{equation}\label{eq:lagrange-observable}
 H(F\compinv(Y))
 =H(Y)+\sum_{n\ge1}\frac{(-1)^n}{n!}
 \D_Y^{\,n-1}\bigl(H'(Y)A(Y)^n\bigr).
\end{equation}
\end{theorem}
\begin{proof}
Use the general coefficient form of Lagrange--B\"urmann in the proof of
\cref{thm:lagrange-transserial}, with $\phi=-A$.
\end{proof}

For example, taking $H(Y)=\log Y$ immediately gives the expansion of the
logarithm of an inverse without first computing and then re-expanding the
inverse itself.

\section{Equivalence with fixed-point iteration}
Formula \eqref{eq:lagrange-main} solves
$G=Y-A(G)$, so uniqueness from \cref{thm:fixed-inverse} shows that it equals
the fixed-point limit.  Conversely, repeatedly substituting
$G=Y-A(G)$ and collecting terms by the number of copies of $A$ reconstructs
the rooted-tree combinatorics encoded by Lagrange inversion.  Newton iteration
converges to the same unique element.  Thus the four methods differ only in
organization:
\begin{center}
\begin{tabularx}{0.95\textwidth}{@{}lX@{}}
\toprule
Method & Organization of the same inverse \\
\midrule
Fixed point & one additional filtration block per substitution \\
Triangular recurrence & one residual coefficient canceled at a time \\
Newton & many blocks corrected simultaneously, with precision doubling \\
Lagrange--B\"urmann & explicit sum indexed by the number of perturbation copies \\
\bottomrule
\end{tabularx}
\end{center}

\section{A direct example}
For $F=X+a\log X$, take $A=aL$.  Then
\begin{align*}
 F\compinv(Y)
 &=Y-aL+a^2Lt
 +a^3\left(\frac12L^2-L\right)t^2\\
 &\quad+a^4\left(\frac13L^3-\frac32L^2+L\right)t^3
 +\FormalBigOAt{t^4}{t}.
\end{align*}
The factor $a^{n+1}$ at block $t^n$ follows from the homogeneity of
$\D^{n}(aL)^{n+1}$.  Setting $a=1$ gives the Wright omega expansion.

\begin{summarybox}
For $F=X+A$, the inverse is
$Y+\sum_{n\ge1}(-1)^n\D^{n-1}(A^n)/n!$.  Differentiation raises outer order,
so the apparently infinite expression is finite at every requested block.
The formula unifies recursive inversion, Bell-polynomial combinatorics, and
the all-orders Lambert expansion.
\end{summarybox}


\chapter{Dominant balance and inverses beyond the identity scale}
\label{ch:reversion-general}

\section{Why dominant balance comes first}
The preceding chapters assume $F(X)=X+\LittleO(X)$.  A general asymptotic inverse may
have a different leading scale.  The correct workflow is:
\begin{enumerate}
\item identify a leading model $F_0$ that can be inverted;
\item write $X=F_0\compinv(Y)$ times or plus a small correction;
\item transform the equation so the remaining unknown is near the identity;
\item apply the polynomial--logarithmic calculus.
\end{enumerate}
Failure to isolate the leading balance can make a formally equivalent fixed
point noncontractive.

\section{Inverting a power--log monomial}
Consider
\begin{equation}\label{eq:power-log-leading}
 Y=cX^\rho(\log X)^\sigma,
 \qquad c>0,
 \qquad \rho>0.
\end{equation}
Put
\[
 S=\log Y,
 \qquad u=\log X.
\]
Then
\begin{equation}\label{eq:power-log-logeq}
 S=\log c+\rho u+\sigma\log u.
\end{equation}
The leading approximation is $u\sim S/\rho$.  Reversion in the variable $u$
gives
\begin{align}\label{eq:u-power-log-inverse}
 u=\frac1\rho\biggl[&S-\sigma\log S+\sigma\log\rho-\log c\\
 &+\frac{\sigma^2\log S+\sigma\log c-\sigma^2\log\rho}{S}
 +\BigOAt{(\log S)^2/S^2}{S\to+\infty}\biggr].\notag
\end{align}
Exponentiating yields
\begin{align}\label{eq:X-power-log-inverse}
 X={}&c^{-1/\rho}\rho^{\sigma/\rho}Y^{1/\rho}
 S^{-\sigma/\rho}\\
 &\times\left[1+
 \frac{\sigma^2\log S+\sigma\log c-\sigma^2\log\rho}
      {\rho S}
 +\BigOAt{(\log S)^2/S^2}{S\to+\infty}\right].\notag
\end{align}
Further terms are generated recursively in powers of $S^{-1}$ with
polynomial coefficients in $\log S$.

\begin{proof}[Derivation of \eqref{eq:u-power-log-inverse}]
Write
\[
 u=\rho^{-1}\left(S-\sigma\log S+C+\frac{a\log S+b}{S}+\cdots\right).
\]
Then
\[
 \log u=\log S-\log\rho
 +\frac{-\sigma\log S+C}{S}
 +\BigOAt{(\log S)^2/S^2}{S\to+\infty}.
\]
Substitution into \eqref{eq:power-log-logeq} gives
$C=\sigma\log\rho-\log c$, $a=\sigma^2$, and
$b=\sigma\log c-\sigma^2\log\rho$.
\end{proof}

\section{An exact affine--logarithmic inverse}
The equation
\begin{equation}\label{eq:affine-log-equation}
 Y=X+a\log X+b,
 \qquad a\ne0,
\end{equation}
has the exact solution
\begin{equation}\label{eq:affine-log-W}
 X=a\,\LambertW_k\!\left(\frac1a
 \exp\!\left(\frac{Y-b}{a}\right)\right),
\end{equation}
with the branch $k$ selected by the desired domain.  Indeed, exponentiating
\eqref{eq:affine-log-equation} after division by $a$ gives
\[
 \frac Xa\exp(X/a)=\frac1a\exp((Y-b)/a).
\]
This identity is both a useful check on formal reversion and a source of
branch information that a purely formal calculation does not record.

\section{Power times exponential}
For
\begin{equation}\label{eq:power-exp-equation}
 X^\alpha\EulerE^{\beta X}=Y,
 \qquad \alpha\beta\ne0,
\end{equation}
one may take an $\alpha$th root and obtain
\[
 X\EulerE^{(\beta/\alpha)X}=Y^{1/\alpha}.
\]
Thus
\begin{equation}\label{eq:power-exp-W}
 X=\frac{\alpha}{\beta}
 \LambertW_k\!\left(\frac{\beta}{\alpha}Y^{1/\alpha}\right).
\end{equation}
The large-$Y$ expansion therefore inherits exactly the
polynomial--logarithmic structure of $\LambertW$ after scaling.

\section{Perturbing a power--log leading term}
Suppose
\begin{equation}\label{eq:perturbed-power-log}
 Y=cX^\rho L^\sigma\bigl(1+E(X)\bigr),
 \qquad E\precdom1.
\end{equation}
Taking logarithms gives
\[
 \log Y=\log c+\rho L+\sigma\log L+\log(1+E).
\]
First invert the unperturbed left side by \eqref{eq:u-power-log-inverse}, then
treat $\log(1+E)$ as a small additive perturbation in the $L$ variable.
This reduces the problem to near-identity reversion.  It is the general
mechanism by which multiplicative asymptotic corrections become additive
polynomial--logarithmic corrections after a logarithm.

\section{Critical points and Puiseux reversal}
If $F'(X_0)=0$, local inversion is not near the identity.  When
\[
 F(X)-F(X_0)=c_m(X-X_0)^m+\cdots,
 \qquad c_m\ne0,
\]
the inverse begins with
\[
 X-X_0\sim c_m^{-1/m}(Y-F(X_0))^{1/m}.
\]
Fractional powers are unavoidable, and there are generally $m$ branches.
Lambert $\LambertW$ near $-\EulerE^{-1}$ is the important case $m=2$: its branch-point
expansion is Puiseux rather than polynomial--logarithmic.  One should not try
to force such a local inverse into the large-argument scale.

\begin{summarybox}
A non-near-identity inverse begins with dominant balance.  Power--log leading
terms are inverted by taking logarithms and reverting an affine-log equation;
exact affine-log and power-exponential equations reduce to Lambert $\LambertW$.
Critical points instead produce Puiseux branches and require a different
local scale.
\end{summarybox}

\part{Lambert \texorpdfstring{$\LambertW$}{W} and the canonical logarithmic inverse}
\label{part:lambert}

\chapter{Wright omega as the inverse of \texorpdfstring{$X+\log X$}{X+log X}}
\label{ch:lambert-omega}

\section{From Lambert \texorpdfstring{$\LambertW$}{W} to an additive inverse problem}
Lambert's function satisfies
\[
 \LambertW(z)\EulerE^{\LambertW(z)}=z.
\]
Use the principal value for the first logarithm and fix compatible analytic
branches $\mathcal B_2$, $\mathcal B_W$, and $\mathcal B_\WrightOmega$ for the
subsequent logarithms.  Put
\begin{equation}\label{eq:L1-L2-def}
 L_{1,k}=\PrincipalLogarithm z+2\pi\ImaginaryUnit k,
 \qquad
 L_{2,k}=\LogarithmOnBranch{\mathcal B_2}{L_{1,k}}.
\end{equation}
Taking logarithms of the defining equation on compatible branches gives
\begin{equation}\label{eq:W-additive-equation}
 \LambertW_k(z)+\LogarithmOnBranch{\mathcal B_W}{\LambertW_k(z)}=L_{1,k}.
\end{equation}
Thus $\LambertW_k(z)$ is obtained by inverting
\[
 F(X)=X+\LogarithmOnBranch{\mathcal B_W}{X}
\]
at the large argument $Y=L_{1,k}$.

The Wright omega function packages this additive problem directly:
\begin{equation}\label{eq:omega-def}
 \WrightOmega(Y)+\LogarithmOnBranch{\mathcal B_\WrightOmega}{\WrightOmega(Y)}=Y,
 \qquad
 \WrightOmega(Y)=\LambertW(\EulerE^Y)
\end{equation}
where branches are chosen consistently.  The DLMF treats both $\LambertW$ and
$\WrightOmega$ and records their branch structure and asymptotic expansions
\cite{DLMF413}.  For formal calculation, \eqref{eq:omega-def} simply says
\begin{equation}\label{eq:omega-inverse}
 \WrightOmega=(X+\log X)\compinv.
\end{equation}

\section{The first terms by fixed-point iteration}
Equation \eqref{eq:omega-def} is
\begin{equation}\label{eq:omega-fixed}
 \WrightOmega=Y-\log\WrightOmega.
\end{equation}
The dominant term is $Y$.  Starting with $\WrightOmega_0=Y$ gives
\[
 \WrightOmega_1=Y-L,
 \qquad L=\log Y.
\]
Substitute again:
\begin{align*}
 \WrightOmega_2
 &=Y-\log(Y-L)\\
 &=Y-L-\log(1-L/Y)\\
 &=Y-L+\frac{L}{Y}+\frac{L^2}{2Y^2}
   +\frac{L^3}{3Y^3}+\cdots.
\end{align*}
This is not yet correct at every displayed block because the correction inside
the logarithm must itself be iterated.  Coefficient stabilization or the
closed formula below produces
\begin{align}\label{eq:omega-first-terms}
 \WrightOmega(Y)={}&Y-L+\frac{L}{Y}
 +\frac{\frac12L^2-L}{Y^2}
 +\frac{\frac13L^3-\frac32L^2+L}{Y^3}\notag\\
 &+\frac{\frac14L^4-\frac{11}{6}L^3+3L^2-L}{Y^4}
 +\BigOAt{L^5/Y^5}{|Y|\to\infty,\ Y\in S_k}.
\end{align}
The first two scales $Y$ and $L=\log Y$ determine the entire hierarchy.

\section{Derivation by Lagrange--B\"urmann}
In \cref{thm:lagrange-transserial}, take $A(Y)=\log Y=L$.  Then
\begin{equation}\label{eq:omega-lagrange}
 \WrightOmega(Y)
 =Y+\sum_{n\ge1}\frac{(-1)^n}{n!}
 \D_Y^{\,n-1}L^n.
\end{equation}
The term indexed by $n$ lies entirely in outer block $Y^{-(n-1)}$.  Therefore
there is exactly one Lagrange term per inverse-power block---an exceptional
simplification that explains the clean all-orders formula.

Writing $r=n-1$ gives
\begin{equation}\label{eq:omega-P-expansion}
 \WrightOmega(Y)=Y-L+\sum_{r\ge1}\frac{\LambertPolynomial{r}(L)}{Y^r},
\end{equation}
where $\LambertPolynomial{r}$ is a polynomial of degree $r$.  The next chapter develops these
polynomials in detail.

\section{A formal generating equation}
Introduce independent formal variables $s$ and $z$, and define
\begin{equation}\label{eq:Q-generating}
 Q(s,z)=\sum_{r\ge1}\LambertPolynomial{r}(z)s^r.
\end{equation}
Substitute
\[
 \WrightOmega(Y)=Y-z+Q(s,z),
 \qquad s=Y^{-1},\quad z=\log Y,
\]
into $\WrightOmega+\log\WrightOmega=Y$.  Since
\[
 \log\WrightOmega=z+\log(1-sz+sQ),
\]
one obtains the compact implicit equation
\begin{equation}\label{eq:Q-implicit}
 Q=-\log(1-sz+sQ).
\end{equation}
Conversely, the unique solution of \eqref{eq:Q-implicit} with $Q(0,z)=0$
generates all Lambert polynomials.  This is a two-variable ordinary power
series problem whose coefficient ring is $\K[z]$.

\section{A numerical illustration}
Let $Y=100$, so the corresponding Lambert argument is $z=\EulerE^{100}$.  The exact
value is approximately
\[
 \LambertW(\EulerE^{100})=95.44148664557583184017204\ldots.
\]
Successive truncations of \eqref{eq:omega-P-expansion} have absolute errors:
\begin{center}
\begin{tabular}{@{}c c@{}}
\toprule
Largest retained block & Absolute error \\
\midrule
$Y-L$ & $4.67\times10^{-2}$ \\
$\LambertPolynomial{1}/Y$ & $6.05\times10^{-4}$ \\
$\LambertPolynomial{2}/Y^2$ & $5.27\times10^{-6}$ \\
$\LambertPolynomial{3}/Y^3$ & $8.15\times10^{-8}$ \\
$\LambertPolynomial{4}/Y^4$ & $5.60\times10^{-9}$ \\
$\LambertPolynomial{5}/Y^5$ & $1.54\times10^{-10}$ \\
$\LambertPolynomial{6}/Y^6$ & $2.03\times10^{-12}$ \\
\bottomrule
\end{tabular}
\end{center}
The table is illustrative rather than an error theorem.  A residual-based
bound is developed in \cref{ch:lambert-branches}.

\section{Why the Lambert expansion is the model example}
The expansion displays almost every structural feature of the present
calculus in its simplest form:
\begin{itemize}
\item a logarithm appears because the defining equation contains an
      exponential;
\item a second logarithm appears when the leading term is substituted back;
\item each inverse power of the first logarithm carries a polynomial in the
      second;
\item composition and reversion are essential, not optional notation;
\item complex branches are represented by changing $L_{1,k}$;
\item a different scale, namely a square-root Puiseux series, takes over at
      the branch point $-\EulerE^{-1}$.
\end{itemize}

\begin{summarybox}
The Wright omega function is the compositional inverse of $X+\log X$.
Lambert $\LambertW_k(z)$ is obtained by evaluating this inverse at
$L_{1,k}=\PrincipalLogarithm z+2\pi\ImaginaryUnit k$.  Lagrange inversion contributes exactly one term
to each inverse-power block, producing polynomials in
$L_{2,k}=\LogarithmOnBranch{\mathcal B_2}{L_{1,k}}$.
\end{summarybox}


\chapter{Lambert polynomials and Stirling numbers}
\label{ch:lambert-polynomials}

\section{The differential-operator formula}
By \cref{thm:repeated-derivative},
\[
 \D_Y^r L^{r+1}
 =Y^{-r}\prod_{j=0}^{r-1}(\partial_L-j)L^{r+1}.
\]
Comparing with \eqref{eq:omega-lagrange} gives the first all-orders formula.

\begin{theorem}[Differential formula for the Lambert polynomials]
\label{thm:P-differential}
For $r\ge1$,
\begin{equation}\label{eq:P-differential}
 \LambertPolynomial{r}(z)=\frac{(-1)^{r+1}}{(r+1)!}
 \prod_{j=0}^{r-1}(\partial_z-j)z^{r+1}.
\end{equation}
\end{theorem}
The shared family also includes the index-zero term
$\LambertPolynomial{0}(z)=-z$.  Consequently
\[
 \WrightOmega(Y)=Y+
 \sum_{r\ge0}\frac{\LambertPolynomial{r}(L)}{Y^r},
 \qquad L=\log Y,
\]
which is identical to the convention that displays $-L$ separately and sums
over $r\ge1$.
The product is a falling factorial of the differential operator,
$\partial_z^{\underline r}$.

\section{The Stirling-number formula}
Let $\stirlingone{r}{q}$ denote the unsigned Stirling number of the first kind,
the number of permutations of $r$ elements with $q$ cycles.  The signed
falling-factorial identity is
\begin{equation}\label{eq:stirling-operator}
 \prod_{j=0}^{r-1}(D-j)
 =\sum_{q=0}^{r}(-1)^{r-q}\stirlingone{r}{q}D^q.
\end{equation}
Applying this to $z^{r+1}$ in \eqref{eq:P-differential} yields:

\begin{theorem}[Stirling formula]\label{thm:P-stirling}
For every $r\ge1$,
\begin{equation}\label{eq:P-stirling}
 \LambertPolynomial{r}(z)=\sum_{m=1}^{r}
 \frac{(-1)^{r-m}}{m!}
 \stirlingone{r}{r-m+1}z^m.
\end{equation}
Consequently,
\begin{equation}\label{eq:W-double-series}
 \LambertW_k(z)\sim L_{1,k}-L_{2,k}
 +\sum_{r\ge1}\sum_{m=1}^{r}
 \frac{(-1)^{r-m}}{m!}
 \stirlingone{r}{r-m+1}
 \frac{L_{2,k}^m}{L_{1,k}^r}.
\end{equation}
\end{theorem}
\begin{proof}
In \eqref{eq:stirling-operator}, apply $D^q$ to $z^{r+1}$ and put
$m=r+1-q$.  The factorial $(r+1)!/m!$ cancels the denominator in
\eqref{eq:P-differential}; the signs reduce to $(-1)^{r-m}$.
\end{proof}

Formula \eqref{eq:W-double-series} is the standard Stirling-number expansion
recorded by Corless et al. and the DLMF \cite{Corless1996,DLMF413}.

\section{A differential and integral recurrence}
The Stirling recurrence
\[
 \stirlingone{r+1}{q}=r\stirlingone{r}{q}
 +\stirlingone{r}{q-1}
\]
translates into a simple recurrence for $\LambertPolynomial{r}$.

\begin{theorem}[Lambert-polynomial recurrence]\label{thm:P-recurrence}
With $\LambertPolynomial{1}(z)=z$,
\begin{equation}\label{eq:P-differential-recurrence}
 \LambertPolynomial{r+1}'(z)=r\LambertPolynomial{r}(z)-\LambertPolynomial{r}'(z),
 \qquad \LambertPolynomial{r+1}(0)=0.
\end{equation}
Equivalently,
\begin{equation}\label{eq:P-integral-recurrence}
 \LambertPolynomial{r+1}(z)=r\int_0^z\LambertPolynomial{r}(u)\Differential u-\LambertPolynomial{r}(z).
\end{equation}
\end{theorem}
This recurrence computes one polynomial using one integration and one
subtraction, with exact rational coefficients.

If $\LambertPolynomial{r}(z)=\sum_m p_{r,m}z^m$, then
\begin{equation}\label{eq:P-coefficient-recurrence}
 p_{r+1,m+1}=\frac{r}{m+1}p_{r,m}-p_{r,m+1},
\end{equation}
where missing coefficients are zero.  This is convenient for an array-based
implementation.

\section{The first ten polynomials}
The first eight fit comfortably on the page:
\begin{align*}
 \LambertPolynomial{1}(z)={}&z,\\
 \LambertPolynomial{2}(z)={}&\frac12z^2-z,\\
 \LambertPolynomial{3}(z)={}&\frac13z^3-\frac32z^2+z,\\
 \LambertPolynomial{4}(z)={}&\frac14z^4-\frac{11}{6}z^3+3z^2-z,\\
 \LambertPolynomial{5}(z)={}&\frac15z^5-\frac{25}{12}z^4
              +\frac{35}{6}z^3-5z^2+z,\\
 \LambertPolynomial{6}(z)={}&\frac16z^6-\frac{137}{60}z^5
              +\frac{75}{8}z^4-\frac{85}{6}z^3
              +\frac{15}{2}z^2-z,\\
 \LambertPolynomial{7}(z)={}&\frac17z^7-\frac{49}{20}z^6
              +\frac{203}{15}z^5-\frac{245}{8}z^4
              +\frac{175}{6}z^3-\frac{21}{2}z^2+z,\\
 \LambertPolynomial{8}(z)={}&\frac18z^8-\frac{363}{140}z^7
              +\frac{3283}{180}z^6-\frac{6769}{120}z^5
              +\frac{245}{3}z^4-\frac{161}{3}z^3
              +14z^2-z.
\end{align*}
For reference, the next two are
\begin{align*}
 \LambertPolynomial{9}(z)={}&\frac19z^9-\frac{761}{280}z^8
 +\frac{29531}{1260}z^7-\frac{1869}{20}z^6
 +\frac{7483}{40}z^5\\
 &-189z^4+91z^3-18z^2+z,\\
 \LambertPolynomial{10}(z)={}&\frac1{10}z^{10}-\frac{7129}{2520}z^9
 +\frac{6515}{224}z^8-\frac{9046}{63}z^7
 +\frac{5985}{16}z^6\\
 &-\frac{21091}{40}z^5+\frac{1575}{4}z^4
 -145z^3+\frac{45}{2}z^2-z.
\end{align*}

\section{Structural properties}
The formulas immediately imply:
\begin{enumerate}
\item $\LambertPolynomial{r}(0)=0$.
\item $\deg \LambertPolynomial{r}=r$ and the leading coefficient is $1/r$.
\item The linear coefficient is $(-1)^{r+1}$.
\item The nonzero coefficients alternate in sign when read from highest to
      lowest degree.
\item All coefficients are rational, and multiplying the coefficient of
      $z^m$ by $m!$ gives an integer Stirling number up to sign.
\end{enumerate}

The highest-degree term alone gives
\[
 \frac{\LambertPolynomial{r}(L_2)}{L_1^r}
 =\frac1r\left(\frac{L_2}{L_1}\right)^r
 +\hbox{lower powers of }L_2.
\]
This explains why the natural small ratio is $L_2/L_1$.

\section{Independent coefficient checks}
A high-order computation has several inexpensive checks:
\begin{enumerate}
\item verify the recurrence \eqref{eq:P-differential-recurrence};
\item verify $\LambertPolynomial{r}(0)=0$, degree $r$, and leading coefficient $1/r$;
\item compare coefficients with \eqref{eq:P-stirling};
\item substitute the truncated $\LambertW$ expansion into $\LambertW+\log \LambertW=L_1$ and check
      that the residual begins one block later.
\end{enumerate}
Agreement of these structurally different tests is much stronger than a
single numerical comparison.

\begin{summarybox}
The coefficient of $L_1^{-r}$ in Lambert's large-argument expansion is a
degree-$r$ polynomial $\LambertPolynomial{r}(L_2)$.  It has equivalent differential-operator,
Stirling-number, differential-recurrence, and integral-recurrence
representations.  These formulas give exact all-orders generation and strong
internal consistency checks.
\end{summarybox}


\chapter{Branches, convergence, and error control for Lambert \texorpdfstring{$\LambertW$}{W}}
\label{ch:lambert-branches}

\section{The branch-indexed expansion}
For a fixed integer $k$, set $L_{1,k}$ and $L_{2,k}$ as in
\eqref{eq:L1-L2-def}, with compatible logarithm branches.  Then, as
$\AbsoluteValue{L_{1,k}}\to\infty$ in an admissible sector,
\begin{equation}\label{eq:Wk-expansion}
 \LambertW_k(z)\sim L_{1,k}-L_{2,k}
 +\sum_{r\ge1}\frac{\LambertPolynomial{r}(L_{2,k})}{L_{1,k}^r}.
\end{equation}
The branch number enters first through the additive term
$2\pi\ImaginaryUnit k$ in
$L_{1,k}$; the rest of the coefficient machinery is universal.  Corless et
al. give the classical branch analysis and all-branch expansion
\cite{Corless1996}; the DLMF records the current standard notation and branch
cuts \cite{DLMF413}.

\section{The principal real branch at infinity}
For $x>0$, take
\[
 L_1=\log x,
 \qquad
 L_2=\log\log x.
\]
Then
\begin{align}\label{eq:W0-real-expansion}
 \LambertW_0(x)={}&L_1-L_2+\frac{L_2}{L_1}
 +\frac{L_2^2-2L_2}{2L_1^2}\\
 &+\frac{2L_2^3-9L_2^2+6L_2}{6L_1^3}
 +\frac{3L_2^4-22L_2^3+36L_2^2-12L_2}{12L_1^4}
 +\cdots.\notag
\end{align}
Every truncation is real once $x>1$ and $L_2$ is defined.  For numerical work
not very far from the branch point, however, a local series or Newton method
is usually preferable to a large-argument expansion.

\section{The lower real branch near zero}
Let $0<\LowerBranchParameter\to0$ and write
\[
 \LambertW_{-1}(-\LowerBranchParameter)=-u,
 \qquad u>0.
\]
The equation $\LambertW\EulerE^\LambertW=-\LowerBranchParameter$ becomes
\[
 u\EulerE^{-u}=\LowerBranchParameter.
\]
With
\[
 S=\log(1/\LowerBranchParameter),
 \qquad L=\log S,
\]
we have
\begin{equation}\label{eq:lower-branch-inverse}
 u-\log u=S.
\end{equation}
This is the inverse of $X-\log X$.  Applying
\cref{thm:lagrange-transserial} with $A=-\log X$ gives
\begin{equation}\label{eq:u-lower-branch}
 u=S+L+\sum_{r\ge1}
 \frac{(-1)^{r+1}\LambertPolynomial{r}(L)}{S^r}.
\end{equation}
Therefore
\begin{align}\label{eq:Wminus1-expansion}
 \LambertW_{-1}(-\LowerBranchParameter)={}&-S-L-\frac{L}{S}
 -\frac{L-\frac12L^2}{S^2}\\
 &-\frac{L-\frac32L^2+\frac13L^3}{S^3}
 -\frac{L-3L^2+\frac{11}{6}L^3-\frac14L^4}{S^4}
 +\cdots.\notag
\end{align}
The same Lambert polynomials reappear with alternating block signs.

\section{The branch point uses a different scale}
At $z=-\EulerE^{-1}$, the derivative of $w\mapsto w\EulerE^w$ vanishes at $w=-1$.
Put
\[
 p=\sqrt{2(1+\EulerE z)}.
\]
On the principal side,
\begin{equation}\label{eq:W-branchpoint}
 \LambertW_0(z)=-1+p-\frac13p^2+\frac{11}{72}p^3
 -\frac{43}{540}p^4+\frac{769}{17280}p^5+\cdots,
\end{equation}
while the adjacent branch changes the signs of the odd powers.  This is a
Puiseux expansion in a square root.  It is not a failure of
polynomial--logarithmic asymptotics; it is a different local geometry caused
by a critical point.

\section{Formal asymptotics versus convergence}
Equation \eqref{eq:Q-implicit} and the analytic implicit-function theorem show
that, for each fixed $z$, the series $Q(s,z)$ converges for sufficiently small
$\AbsoluteValue{s}$.  Along the Lambert path $s=1/L_1$, $z=L_2$, the quantity
$\AbsoluteValue{s z}=\AbsoluteValue{L_2/L_1}$ tends to zero.  Hence the series is genuinely convergent in
a sufficiently far large-argument region on a fixed branch.  Outside a
convergence region, its finite truncations still provide the standard
Poincar\'e asymptotic expansion in admissible sectors.  Precise domains and
branch qualifications are given in the classical Lambert-$\LambertW$ literature
\cite{Corless1996,DLMF413}.

\section{Residual-based error estimation}
Let
\begin{equation}\label{eq:W-truncation}
 w_N=L_1-L_2+\sum_{r=1}^{N}\frac{\LambertPolynomial{r}(L_2)}{L_1^r}
\end{equation}
and define the additive residual
\begin{equation}\label{eq:W-residual}
 R_N=w_N+\LogarithmOnBranch{\mathcal B_W}{w_N}-L_1.
\end{equation}
The exact solution $w$ satisfies $F(w)=L_1$ for
$F(u)=u+\LogarithmOnBranch{\mathcal B_W}{u}$.  Since
$F'(u)=1+1/u$, one Newton correction is
\begin{equation}\label{eq:W-residual-correction}
 \delta_N=-\frac{R_N}{1+1/w_N}.
\end{equation}
Taylor's theorem gives
\begin{equation}\label{eq:W-error-from-residual}
 w-w_N=\delta_N+
 \BigOAt{R_N^2/\AbsoluteValue{w_N}^2}
        {|L_1|\to\infty,\ L_1\in S_k}
\end{equation}
provided the segment or complex neighborhood stays away from $0$ and the
chosen branch cut.  Thus the residual supplies both a check and an improved
approximation.

For the corresponding truncation on a fixed branch sector,
\[
 R_N=\BigOAt{L_2^{N+1}/L_1^{N+1}}
              {|L_1|\to\infty,\ L_1\in S_k},
\]
so the inverse error has the same leading order because $F'(w)=1+\LittleO(1)$.
This is the analytic reflection of the triangular residual-cancellation
algorithm.

\section{Branch bookkeeping rules}
In complex calculations record all of the following:
\begin{enumerate}
\item the principal value $\PrincipalLogarithm z$ used to form $L_{1,k}$;
\item the integer $k$;
\item the branch $\mathcal B_2$ used in
      $\LogarithmOnBranch{\mathcal B_2}{L_{1,k}}=L_{2,k}$;
\item the compatible branches $\mathcal B_W$ and $\mathcal B_\WrightOmega$ used
      for the additive inverse equations;
\item the sector in which $\AbsoluteValue{L_{1,k}}\to\infty$;
\item any continuation path crossing or avoiding branch cuts.
\end{enumerate}
Two algebraically identical-looking expansions can represent different
branches if these data differ.

\begin{summarybox}
The same Lambert polynomials govern every large complex branch through
$L_{1,k}=\PrincipalLogarithm z+2\pi\ImaginaryUnit k$.  The lower real branch near zero is the inverse of
$X-\log X$ and uses alternating polynomial blocks.  Near $-\EulerE^{-1}$ the scale
changes to a square-root Puiseux series.  Residual evaluation turns a formal
truncation into a practical error estimator and Newton correction.
\end{summarybox}


\chapter{Generalized Lambert-type inversions and applications}
\label{ch:lambert-applications}

\section{Equations reducible exactly to \texorpdfstring{$\LambertW$}{W}}
A surprising number of inverse problems reduce to $\LambertW$ after elementary
normalization.

\subsection{Power times exponential}
As shown in \eqref{eq:power-exp-W},
\[
 X^\alpha\EulerE^{\beta X}=Y
 \quad\Longrightarrow\quad
 X=\frac{\alpha}{\beta}
 \LambertW_k\!\left(\frac{\beta}{\alpha}Y^{1/\alpha}\right).
\]
Its large-$Y$ expansion follows by scaling $L_1$ and $L_2$.

\subsection{A linear term plus a logarithm}
Equation $Y=X+a\log X+b$ has the exact inverse
\[
 X=aW_k\!\left(a^{-1}\EulerE^{(Y-b)/a}\right).
\]
Expanding this exact expression reproduces the direct reversion coefficients
from \cref{ch:reversion-triangular}.

\subsection{The inverse of \texorpdfstring{$X\log X$}{X log X}}
If
\[
 Y=X\log X,
\]
put $u=\log X$.  Then $u\EulerE^u=Y$, so
\begin{equation}\label{eq:XlogX-inverse}
 X=\EulerE^{\LambertW(Y)}=\frac{Y}{\LambertW(Y)}.
\end{equation}
With $L_1=\log Y$ and $L_2=\log L_1$, division of the Lambert expansion gives
\begin{align}\label{eq:XlogX-asymptotic}
 X=\frac{Y}{L_1}\biggl[1
 &+\frac{L_2}{L_1}
 +\frac{L_2^2-L_2}{L_1^2}\\
 &+\frac{L_2^3-\frac52L_2^2+L_2}{L_1^3}
 +\BigOAt{L_2^4/L_1^4}{|L_1|\to\infty,\ L_1\in S_k}\biggr].\notag
\end{align}
This is a compact example in which reversion is followed by formal division.

\section{Perturbed affine-log equations}
Consider
\begin{equation}\label{eq:perturbed-affine-log}
 Y=X+a\log X+b+\frac{c\log X+d}{X}
 +\frac{e(\log X)^2+f\log X+g}{X^2}+\cdots.
\end{equation}
One may either:
\begin{enumerate}
\item use the exact affine-log inverse as a new leading variable, then compose
      the perturbation around it; or
\item regard the entire right side as $X+A(X)$ and apply direct near-identity
      reversion.
\end{enumerate}
The first approach can reduce intermediate coefficient size when $a$ is not
small; the second is simpler and produces the universal formulas
\eqref{eq:param-g0}--\eqref{eq:param-g2} immediately.

\section{Normal tail and inverse-error-function scales}
Let $p\to0^+$ and let $q$ be the upper standard-normal quantile,
$\Probability(Z>q)=p$.  The classical tail expansion has the form
\[
 p=\frac{\EulerE^{-q^2/2}}{\sqrt{2\pi}\,q}
 \left(1-q^{-2}+3q^{-4}-15q^{-6}+\cdots\right).
\]
Taking logarithms and putting $S=\log(1/p)$ turns the inversion into
\[
 2S=q^2+2\log q+\log(2\pi)+\hbox{inverse powers of }q^2.
\]
With $Y=q^2$, this is an affine-log equation plus a
polynomial--logarithmic perturbation.  The first result is
\begin{equation}\label{eq:normal-quantile-first}
 q=\sqrt{2S}
 -\frac{\log S+\log(4\pi)}{2\sqrt{2S}}
 +\BigOAt{(\log S)^2/S^{3/2}}{S\to+\infty}.
\end{equation}
All later terms can be obtained by the same reversion algorithms.  This
explains why logarithms of logarithms occur in extreme quantiles.

\section{Formal inversion of the logarithmic integral}
The logarithmic integral has the large-$x$ expansion
\begin{equation}\label{eq:li-expansion}
 \operatorname{li}(x)
 \sim\frac{x}{\log x}
 \left(1+\frac1{\log x}+\frac{2!}{(\log x)^2}
 +\frac{3!}{(\log x)^3}+\cdots\right).
\end{equation}
Its inverse therefore belongs to the same power--logarithmic reversion class.
Writing $L=\log n$ and $\ell=\log L$, formal inversion gives the Cipolla-type
expansion
\begin{align}\label{eq:nth-prime-formal}
 x\sim n\biggl[&L+\ell-1+\frac{\ell-2}{L}
 -\frac{\ell^2-6\ell+11}{2L^2}\\
 &+\frac{\ell^3-9\ell^2+23\ell-11}{6L^3}
 +\cdots\biggr].\notag
\end{align}
When $x$ is specialized to the $n$th prime, analytic number theory is needed
to transfer the formal inverse of $\operatorname{li}$ to $p_n$; the algebraic
coefficient calculation itself is precisely polynomial--logarithmic series
reversal.  This example shows that the calculus is not specific to explicitly
Lambert-solvable equations.

\section{Stirling inversion}
Suppose $N$ is large and one wants to solve
$\GammaFunction(X+1)=N$.  Stirling's
formula gives
\[
 \log N=X\log X-X+\frac12\log(2\pi X)
 +\frac1{12X}+\cdots.
\]
The dominant inverse is controlled by $X\log X$ and hence by Lambert $\LambertW$;
the remaining terms are polynomial--logarithmic perturbations.  A robust
calculation first uses \eqref{eq:XlogX-inverse} for the leading approximation,
then applies Newton or triangular reversion to the Stirling correction.

\section{A reusable reduction pattern}
Many applications fit the sequence
\[
 \begin{gathered}
 \text{implicit equation}
 \longrightarrow \text{take logarithms}
 \longrightarrow \text{isolate a dominant affine-log map}\\
 \longrightarrow \text{near-identity reversion}.
 \end{gathered}
\]
The technique applies to combinatorial saddle points, delay equations,
large-deviation thresholds, inverse special functions, and algorithmic
complexity equations such as $n\log n=N$.  The details differ, but the same
four operations recur: multiplication, division, composition, and reversal.

\begin{summarybox}
Lambert $\LambertW$ is both a special function and a normal form for asymptotic
inversion.  Exact reductions handle powers times exponentials, affine-log
maps, and $X\log X$.  More complicated inverse problems become perturbations
of these models after taking logarithms, producing the same polynomial blocks
in nested logarithms.
\end{summarybox}


\part{Worked calculations, implementation, and analytic meaning}
\label{part:practice}

\chapter{Worked calculations from beginning to end}
\label{ch:worked-calculations}

The preceding chapters separated the operations in order to prove them
cleanly. Actual calculations mix them. This chapter therefore carries out
several complete examples, including all truncation decisions and residual
checks. Throughout,
\[
 t=Y^{-1},\qquad L=\log Y,
\]
unless another variable is explicitly introduced.

\section{A product and a quotient in logarithmic blocks}
Consider
\begin{align*}
 A&=1+(L+1)t+(L^2-L)t^2+\FormalBigOAt{t^3}{t},\\
 B&=1+(2L-1)t+(3-L)t^2+\FormalBigOAt{t^3}{t}.
\end{align*}
The product is obtained by block convolution:
\begin{align*}
 AB
 &=1+\bigl[(L+1)+(2L-1)\bigr]t\\
 &\quad+\bigl[(L^2-L)+(L+1)(2L-1)+(3-L)\bigr]t^2
   +\FormalBigOAt{t^3}{t}\\
 &=1+3Lt+(3L^2-L+2)t^2+\FormalBigOAt{t^3}{t}.
\end{align*}
For the quotient, write
\[
 \frac AB=1+Q_1t+Q_2t^2+\FormalBigOAt{t^3}{t}.
\]
The triangular division rule gives
\[
 Q_1=(L+1)-(2L-1)=2-L
\]
and
\begin{align*}
 Q_2
 &=(L^2-L)-(3-L)-(2L-1)Q_1\\
 &=3L^2-5L-1.
\end{align*}
Consequently
\begin{equation}\label{eq:worked-quotient}
 \frac AB
 =1+(2-L)t+(3L^2-5L-1)t^2+\FormalBigOAt{t^3}{t}.
\end{equation}
A reliable check is multiplication by the divisor. Substituting
\eqref{eq:worked-quotient} into $B(A/B)$ recovers $A$ through $t^2$. This
simple residual check is worth automating: it detects sign errors, omitted
cross terms, and inadequate internal precision.

\section{A fully mixed polynomial--logarithmic reversion}
Let
\begin{equation}\label{eq:worked-F}
 F(X)=X+2\log X+1
 +\frac{3\log X-1}{X}
 +\frac{(\log X)^2+2}{X^2}.
\end{equation}
We seek $G=F\compinv$ as $Y\to+\infty$. Put
\[
 A(X)=F(X)-X.
\]
The Lagrange--B\"urmann operator formula from
\cref{ch:reversion-lagrange} reads
\begin{equation}\label{eq:worked-LB}
 G(Y)=Y-A+\frac12\D(A^2)-\frac16\D^2(A^3)
      +\frac1{24}\D^3(A^4)+\cdots ,
\end{equation}
where every occurrence of $A$ on the right is evaluated at $Y$ and
\[
 \D=t\frac{\partial}{\partial L}
       -t^2\frac{\partial}{\partial t}.
\]
Only the first four displayed terms can contribute through $t^3$. Indeed,
$A=\BigOAt{L}{Y\to+\infty}$ and each application of $\D$ raises the outer $t$-order by one.
A direct calculation gives
\begin{align}
 G(Y)=Y&-2L-1+\frac{L+3}{Y}
 +\frac{-3L^2+7L-3}{Y^2}\notag\\
 &+\frac{-\frac{32}{3}L^3+25L^2-6L-\frac{59}{6}}{Y^3}
 +\BigOAt{L^4/Y^4}{Y\to+\infty}.
 \label{eq:worked-inverse}
\end{align}

\subsection{How the coefficient blocks arise}
For readers checking the calculation by hand, the contributions in
\eqref{eq:worked-LB}, truncated after $t^3$, are
\begin{align*}
 -A={}&-2L-1+(1-3L)t-(L^2+2)t^2,\\
 \frac12\D(A^2)={}&(4L+2)t+(-6L^2+11L+2)t^2\\
 &+(-4L^3-5L^2+9L-4)t^3,\\
 -\frac16\D^2(A^3)={}&(4L^2-4L-3)t^2\\
 &+(-12L^3+46L^2-11L-\tfrac{17}{2})t^3,\\
 \frac1{24}\D^3(A^4)={}&
 (\tfrac{16}{3}L^3-16L^2-4L+\tfrac83)t^3.
\end{align*}
Their sum is \eqref{eq:worked-inverse}. Notice that a term which is
``higher'' in the Lagrange index can still contribute to the same outer
$t$-block. The filtration, not the number of displayed summands, determines
the stopping rule.

\subsection{Residual verification by generator substitution}
Set $B(Y)=G(Y)-Y$ and $U=tB$. Since
\[
 G(Y)=Y(1+U),
\]
composition into \eqref{eq:worked-F} uses the exact substitutions
\[
 \frac1{G(Y)}=\frac{t}{1+U},
 \qquad
 \log G(Y)=L+\log(1+U).
\]
Expanding these two expressions through $t^4$ and substituting into
$F(G(Y))-Y$ gives
\[
 F(G(Y))-Y
 =\left(28L^4-\frac{211}{3}L^3-15L^2+91L-\frac{59}{3}\right)t^4
 +\BigOAt{L^5/Y^5}{Y\to+\infty}.
\]
Thus every coefficient claimed in \eqref{eq:worked-inverse} is certified.
The nonzero $t^4$ residual is expected: it is the first block not included in
the inverse.

\begin{computebox}
For an inverse through $Y^{-N}$, compute one extra residual block. A zero
residual through $Y^{-N}$ proves the retained coefficients; the first
nonzero block is also useful for estimating the next correction.
\end{computebox}

\section{A nested-logarithmic inverse}
Let
\[
 M=\log L,\qquad
 F(X)=X+a\log X+b\log\log X,
\]
where $a,b$ are constants. Define
\[
 A=aL+bM,\qquad p=a+\frac bL.
\]
Then
\[
 A'=\frac{p}{Y},\qquad
 A''=-\frac{a+b/L+b/L^2}{Y^2}.
\]
The first three Lagrange--B\"urmann terms yield
\begin{align}
 F\compinv(Y)
 =Y-A+\frac{Ap}{Y}
 +\frac{1}{Y^2}\left[
   -Ap^2+\frac{A^2}{2}
   \left(a+\frac bL+\frac b{L^2}\right)\right]
 +\BigOAt{L^4/Y^3}{Y\to+\infty}.
 \label{eq:nested-log-inverse}
\end{align}
This formula illustrates an important closure point. The coefficient of
$Y^{-2}$ is no longer a polynomial in $L$ and $M$ alone: negative powers of
$L$ appear after differentiation. Hence nested-logarithmic reversion
naturally lives in an iterated Laurent coefficient field such as
\[
 \K((L^{-1}))[M],
\]
or in a larger recursively defined transseries field if still deeper
logarithms occur.

\section{Reversion by Newton doubling}
The same inverse \eqref{eq:worked-inverse} can be obtained with far fewer
high-order compositions when many blocks are required. Start with
\[
 G_0(Y)=Y-2L-1,
\]
and perform
\begin{equation}\label{eq:worked-newton}
 G_{j+1}
 =G_j-\frac{F(G_j)-Y}{F'(G_j)}.
\end{equation}
If $F(G_j)-Y=\FormalBigOAt{t^m}{t}$ and
$F'(G_j)=1+\FormalBigOAt{t}{t}$, the next residual is typically
$\FormalBigOAt{t^{2m}}{t}$. Thus the working precision can be doubled at
each stage:
\[
 t,\quad t^2,\quad t^4,\quad t^8,\quad\ldots.
\]
At every stage all products, quotients, logarithms, and compositions are
truncated immediately to the current target. Retaining more terms than the
target only increases intermediate expression swell.

\section{A scale-changing inversion with a Lambert core}
Consider
\begin{equation}\label{eq:worked-scale-change}
 Y=X^\rho(\log X)^\sigma
 \left(1+\frac{c_1}{\log X}
       +\frac{c_2}{(\log X)^2}
       +\BigOAt{(\log X)^{-3}}{X\to+\infty}\right),
 \qquad \rho>0.
\end{equation}
Put $Z=\log X$ and $S=\log Y$. Taking logarithms gives
\begin{align}
 S
 &=\rho Z+\sigma\log Z
   +\frac{c_1}{Z}
   +\frac{c_2-\frac12c_1^2}{Z^2}
   +\BigOAt{Z^{-3}}{Z\to+\infty}.
 \label{eq:worked-log-reduction}
\end{align}
The unperturbed equation $S=\rho Z+\sigma\log Z$ has the exact inverse
\[
 Z_0=
 \begin{cases}
 \displaystyle
 \frac{\sigma}{\rho}
 \LambertW\!\left(\frac{\rho}{\sigma}
          \exp\!\left(\frac S\sigma\right)\right),
 &\sigma\ne0,\\[1.2em]
 S/\rho,&\sigma=0.
 \end{cases}
\]
One may now either revert \eqref{eq:worked-log-reduction} directly in the
large variable $S$, or write $Z=Z_0+\Delta$ and revert only the inverse-power
perturbation. Finally,
\[
 X=\exp Z.
\]
This ``logarithm, normalize, Lambert, perturb'' workflow avoids forcing a
scale-changing problem into a near-identity ansatz too early.

\section{An implicit equation not exactly Lambert-solvable}
Suppose
\begin{equation}\label{eq:worked-implicit}
 Y=X+\alpha\log X+\beta
 +\frac{\gamma\log X+\delta}{X}
 +\BigOAt{(\log X)^2/X^2}{X\to+\infty}.
\end{equation}
The triangular formulas of \cref{ch:reversion-triangular} immediately give
\begin{align}
 X=Y&-\alpha L-\beta\\
 &+\frac{(\alpha^2-\gamma)L+\alpha\beta-\delta}{Y}
 +\BigOAt{L^2/Y^2}{Y\to+\infty}.
 \label{eq:worked-implicit-inverse}
\end{align}
This two-line computation is often all that an application needs. The
important point is not the particular formula but the workflow:
\begin{enumerate}
\item choose the dominant variable so that the map is tangent to the
      identity;
\item write one polynomial or iterated-Laurent coefficient per outer power;
\item solve coefficient blocks in increasing outer order;
\item compose back and inspect the residual.
\end{enumerate}

\chapter{Finite-order algorithms and data structures}
\label{ch:implementation}

Formal existence theorems become useful only when accompanied by
coefficientwise algorithms. Algorithmic inversion of exp--log functions by
multiseries is developed systematically by Salvy and Shackell
\cite{SalvyShackell1999}. The algorithms in this chapter are deliberately
finite: they retain a prescribed principal part together with a finite
representative modulo the additive tail filtration.  This is not a quotient
of the Laurent ring by an ideal: \(t\) is a unit in \(\K[L]((t))\), so the
tail \(t^{N+1}\K[L][[t]]\) is not an ideal of that Laurent ring.  No machine is
asked to store an infinite object.

\section{Sparse block representation}
A truncated one-logarithm transseries
\[
 T=\sum_{n=n_0}^{N}p_n(L)t^n+\FormalBigOAt{t^{N+1}}{t}
\]
can be stored as a sparse map
\[
 n\longmapsto p_n.
\]
The outer key $n$ is a rational number if fractional powers are allowed.
The value $p_n$ is a polynomial in $L$ for $\PLpoly$, a Laurent series in
$L^{-1}$ for $\PLfield$, or another sparse transseries object for nested
logarithms.

Canonicalization consists of three operations:
\begin{enumerate}
\item combine equal outer exponents;
\item remove zero coefficient blocks;
\item normalize every coefficient object recursively.
\end{enumerate}
Exact rational arithmetic should be used whenever the input coefficients are
exact. Floating-point coefficients hide cancellations and make zero testing
unreliable.

\section{A precision object}
A practical implementation needs more than a single integer $N$. A useful
precision descriptor records
\[
 \mathcal P=(n_{\min},n_{\max};d_L;\mathcal P_{\mathrm{inner}}),
\]
where $n_{\max}$ is the largest retained outer exponent, $d_L$ optionally
bounds polynomial degree, and $\mathcal P_{\mathrm{inner}}$ controls an
iterated Laurent coefficient field. The degree bound is not part of the
mathematical definition; it is an implementation safeguard. It should be
used only when a theorem or a posteriori residual test guarantees that no
discarded high logarithmic degree can feed back into a retained block.

\section{Multiplication and reciprocal}
For
\[
 A=\sum A_nt^n,\qquad B=\sum B_nt^n,
\]
the product block is
\[
 (AB)_r=\sum_{i+j=r}A_iB_j.
\]
A sparse implementation iterates only over nonzero pairs. If the supports
have $s_A$ and $s_B$ retained blocks, naive multiplication costs
$\BigO(s_As_B)$ coefficient products; dense consecutive support gives the
usual quadratic convolution.

For $B=b_0+\sum_{n\ge1}b_nt^n$ with invertible $b_0$, its reciprocal
$C=\sum c_nt^n$ satisfies
\begin{equation}\label{eq:implementation-reciprocal}
 c_0=b_0^{-1},\qquad
 c_n=-b_0^{-1}\sum_{j=1}^{n}b_jc_{n-j}.
\end{equation}
The same recurrence works when $b_0$ is a nonconstant invertible element of
the coefficient field.

\begin{algorithm}[H]
\caption{Truncated reciprocal}
\KwIn{$B=\sum_{n=0}^{N}b_nt^n$ with invertible $b_0$}
\KwOut{$C=\trunc{B^{-1}}{N}$}
$c_0\leftarrow b_0^{-1}$\;
\For{$n\leftarrow1$ \KwTo $N$}{
  $s\leftarrow0$\;
  \For{$j\leftarrow1$ \KwTo $n$}{
    $s\leftarrow s+b_jc_{n-j}$\;
  }
  $c_n\leftarrow-b_0^{-1}s$\;
}
\Return $\sum_{n=0}^{N}c_nt^n$\;
\end{algorithm}

\section{Differentiation as an operator on blocks}
The derivative is not obtained by treating $t$ and $L$ as independent:
\[
 \frac{\Differential t}{\Differential X}=-t^2,\qquad
 \frac{\Differential L}{\Differential X}=t.
\]
Consequently
\begin{equation}\label{eq:implementation-D}
 \D=t\partial_L-t^2\partial_t
\end{equation}
and
\begin{equation}\label{eq:implementation-D-block}
 \D\!\left(P(L)t^n\right)
 =\bigl(P'(L)-nP(L)\bigr)t^{n+1}.
\end{equation}
Repeated differentiation can be applied blockwise by
\begin{equation}\label{eq:implementation-Dm}
 \D^m\!\left(P(L)t^n\right)
 =t^{n+m}
 \prod_{j=0}^{m-1}
 \left(\partial_L-(n+j)\right)P(L).
\end{equation}
Formula \eqref{eq:implementation-Dm} is substantially faster than expanding
and differentiating the whole expression at every Lagrange step.

\section{Composition by the two-generator interface}
Let
\[
 G(Y)=Y(1+U),\qquad \nu_tU>0.
\]
Every one-logarithm block series can be composed into $G$ by the interface
\begin{equation}\label{eq:implementation-generator-interface}
 t_X\longmapsto t_Y(1+U)^{-1},\qquad
 L_X\longmapsto L_Y+\log(1+U).
\end{equation}
This separates composition into standard operations:
reciprocal, logarithm, polynomial evaluation, powering, multiplication, and
summation. The order of evaluation matters for efficiency. Compute
$(1+U)^{-1}$ and $\log(1+U)$ once, cache their powers, and reuse them for
every source block.

\begin{algorithm}[H]
\caption{Generator-based composition}
\KwIn{$F=\sum_{n=n_0}^{N}p_n(L)t^n$, $G=Y(1+U)$, target order $N$}
\KwOut{$F\compose G$ through outer $t_Y$-block $N$}
$R\leftarrow(1+U)^{-1}$, truncated through $t_Y$-block $N$\;
$H\leftarrow\log(1+U)$, truncated through $t_Y$-block $N$\;
$S\leftarrow0$\;
\ForEach{retained block $(n,p_n)$ of $F$}{
  $S\leftarrow S+t_Y^n R^n\,p_n(L_Y+H)$\;
  truncate $S$ through $t_Y$-block $N$\;
}
\Return $S$\;
\end{algorithm}

For negative $n$, $R^n=(1+U)^{-n}$ is an ordinary positive power. For
rational $n$, use the binomial series. Terms should be truncated after every
multiplication.

\section{Triangular reversion}
For
\[
 F(X)=X+\sum_{n\ge0}a_n(\log X)X^{-n},
\qquad
 G(Y)=Y+\sum_{n\ge0}g_n(\log Y)Y^{-n},
\]
the coefficient $g_N$ occurs linearly and for the first time in the
$Y^{-N}$ block of $F(G)-Y$. This gives the following robust algorithm.

\begin{algorithm}[H]
\caption{Triangular polynomial--logarithmic reversion}
\KwIn{$F=X+A(X)$, target outer order $N$}
\KwOut{$G=F\compinv\bmod Y^{-N-1}$}
$G\leftarrow Y$\;
\For{$n\leftarrow0$ \KwTo $N$}{
  introduce an unknown coefficient block $g_n(L)Y^{-n}$ in $G$\;
  compute only the $Y^{-n}$ block of $F(G)-Y$\;
  solve that block equation for $g_n(L)$\;
  substitute the solution and canonicalize\;
}
\Return $G$\;
\end{algorithm}

The ``solve'' step is normally just negation because the tangent map is the
identity. If the leading derivative is a nonzero coefficient $c$, divide
by $c$. If the derivative vanishes, ordinary transserial reversion is not
the correct local model; fractional powers or a different dominant balance
are required.

\section{Newton reversion with precision doubling}
For high order, Newton iteration
\[
 G_{\mathrm{new}}
 =G-\frac{F(G)-Y}{F'(G)}
\]
usually outperforms one-block-at-a-time reversal. The implementation should
not compute every iterate at full final precision. If the current result is
correct through order $m$, compute the next one only through order $2m$.
This makes the number of composition stages logarithmic in the requested
order.

A safe schedule is:
\[
 1,\;2,\;4,\;8,\;\ldots,\;2^{\Ceiling{\log_2N}},
\]
with the last stage truncated to $N$. Before each quotient, normalize the
leading block of $F'(G)$ and verify that it is invertible.

\section{Lagrange--B\"urmann as a streaming algorithm}
The formal identity
\[
 G-Y=\sum_{n\ge1}\frac{(-1)^n}{n!}\D^{n-1}(A^n)
\]
does not require construction of all terms at once. Maintain the current
power $A^n$ and apply the block operator \eqref{eq:implementation-Dm}. Stop
when the valuation lower bound
\[
 \nu_t\!\left(\D^{n-1}(A^n)\right)>N
\]
guarantees that no retained block can occur. This is particularly effective
when $A$ has positive outer order. When $A$ begins at order $t^0$ but grows
only logarithmically, the derivative count still forces the $n$th Lagrange
term to begin no earlier than $t^{n-1}$.

\section{A compact symbolic skeleton}
The following pseudocode emphasizes the architecture rather than a
particular computer-algebra system.

\begin{lstlisting}[language=Python,caption={Sparse outer blocks with recursive coefficients}]
class PLSeries:
    # blocks[n] is the coefficient of t**n
    def __init__(self, blocks, precision):
        self.blocks = canonicalize(blocks, precision)
        self.precision = precision

    def truncate(self, N):
        return PLSeries({n: p for n, p in self.blocks.items()
                         if n <= N}, N)

    def mul(self, other, N):
        out = {}
        for i, ai in self.blocks.items():
            for j, bj in other.blocks.items():
                if i + j <= N:
                    out[i+j] = out.get(i+j, 0) + ai*bj
        return PLSeries(out, N)

    def derivative(self, N):
        out = {}
        for n, p in self.blocks.items():
            if n + 1 <= N:
                out[n+1] = diff_L(p) - n*p
        return PLSeries(out, N)

    def compose_near_identity(self, U, N):
        R = reciprocal(one() + U, N)
        H = log_series(one() + U, N)
        out = zero()
        for n, p in self.blocks.items():
            term = t_power(n) * power(R, n, N)
            term *= substitute_L(p, L + H)
            out = (out + term).truncate(N)
        return out

def inverse_newton(F, G0, N):
    G, m = G0, 1
    while m < N:
        target = min(2*m, N)
        residual = compose(F, G, target) - identity()
        jacobian = compose(derivative(F), G, target)
        G = (G - divide(residual, jacobian, target)).truncate(target)
        m = target
    return G
\end{lstlisting}

The recursive operations \texttt{diff\_L}, \texttt{substitute\_L},
and \texttt{canonicalize} select polynomial, Laurent, or deeper
transseries coefficients.

\section{Internal precision and guard blocks}
Composition and division may require source terms beyond the desired output
order if leading exponents are negative or if a scale change occurs.
Accordingly, a production implementation should determine internal
precision from valuations rather than simply requesting the same order from
every operand.

For near-identity one-logarithm composition with nonnegative outer
corrections, one or two guard blocks are usually enough for a residual check,
but ``usually'' is not a theorem. The rigorous rule is dependency tracing:
for each output exponent, enumerate all products and substitutions capable of
reaching it. Ratio-set or grid-based transseries implementations formalize
this bookkeeping by recording the small generators needed for a
calculation; see Edgar's grid formalism and van der Hoeven's computational
framework \cite{EdgarBeginners,vdHoeven2006}.

\section{Complexity and expression swell}
At fixed logarithmic degree, dense multiplication and triangular division
are quadratic in the number of retained outer blocks. Newton inversion
reduces the number of high-cost composition stages, and fast polynomial or
series arithmetic can lower asymptotic complexity further. In practice the
main obstacle is often coefficient swell rather than operation count:
polynomial degrees rise, rational numerators grow, and nested coefficient
fields expand.

The most effective controls are:
\begin{enumerate}
\item truncate after every primitive operation;
\item factor or collect only at block boundaries;
\item cache $\log(1+U)$, $(1+U)^\alpha$, and repeated derivatives;
\item postpone full expansion of coefficient polynomials when possible;
\item verify by a residual rather than repeatedly simplifying two enormous
      expressions to canonical equality.
\end{enumerate}


\chapter{From formal transseries to asymptotic expansions}
\label{ch:analytic-meaning}

A formal transseries is an algebraic object. An asymptotic expansion is a
statement about a function and a limit. The two viewpoints cooperate, but
they are not interchangeable without hypotheses.

\section{Poincar\'e meaning in an ordered scale}
Let $(\phi_\nu)$ be a totally ordered asymptotic scale, so that
$\phi_{\nu+1}=\LittleO(\phi_\nu)$ in the chosen enumeration. A function $f$ has
the expansion
\[
 f\sim\sum_{\nu}c_\nu\phi_\nu
\]
when, after every finite truncation at $\nu=N$,
\[
 f-\sum_{\nu\le N}c_\nu\phi_\nu=\LittleO(\phi_N).
\]
For the polynomial--logarithmic scale, a convenient block formulation is
\begin{equation}\label{eq:block-asymptotic-analytic}
 f(X)-\sum_{n=n_0}^{N}X^{-n}p_n(\log X)
 =o\!\left(X^{-N}(\log X)^{d_N}\right)
\end{equation}
after the monomials inside the last block have themselves been ordered.
One must specify that inner ordering; the symbol $\BigO(X^{-N})$ alone does
not distinguish $X^{-N}\log^{100}X$ from $X^{-N}/\log X$.

\section{Uniqueness of coefficients}
\begin{proposition}[Uniqueness in an asymptotic scale]
If the monomials of a polynomial--logarithmic expansion form a strictly
decreasing asymptotic scale and two expansions represent the same function
to every finite order, then their coefficients agree term by term.
\end{proposition}

\begin{proof}
Subtract the two expansions. If a first nonzero coefficient existed, divide
the difference by its monomial. All later monomials tend to zero relative
to it, so the quotient would tend to that nonzero coefficient. But the
assumed remainder statement says the same quotient tends to zero, a
contradiction.
\end{proof}

This elementary argument is the analytic counterpart of uniqueness of Hahn
or grid-based transseries support.

\section{Transfer under arithmetic}
Finite asymptotic expansions are stable under addition and multiplication.
Division is stable when the divisor has a nonvanishing leading term. More
precisely, suppose
\[
 f=F_N+r_f,\qquad g=G_N+r_g,
\]
where the remainders are smaller than the last retained monomials and the
leading term of $G_N$ dominates $r_g$. Then formal multiplication of $F_N$
and $G_N$, followed by truncation, gives the asymptotic product. Similarly,
the geometric expansion
\[
 \frac1g=\frac1{G_N}\frac1{1+r_g/G_N}
\]
justifies the formal reciprocal to every order for which
$r_g/G_N\to0$.

Differentiation is more delicate: a Poincar\'e expansion cannot always be
differentiated termwise. Sufficient conditions include analyticity with
uniform sectorial remainder bounds, or direct derivative estimates on the
remainders. Formal differentiation is always defined in the transseries
algebra; analytic realization requires this extra regularity.

\section{Composition needs uniform control}
Suppose $f(X)$ has an expansion as $X\to+\infty$ and
$g(Y)\to+\infty$. Substituting $X=g(Y)$ into each displayed term is not
enough. The remainder estimate for $f$ must be uniform over the range
traversed by $g(Y)$, and the perturbation of logarithms and powers must
remain inside the selected branch or sector.

For a real near-identity map
\[
 g(Y)=Y(1+u(Y)),\qquad u(Y)\to0,
\]
these conditions are usually straightforward:
\[
 \log g=\log Y+\log(1+u),\qquad
 g^{-\alpha}=Y^{-\alpha}(1+u)^{-\alpha}.
\]
For complex composition, one must additionally keep $1+u$ away from the
chosen logarithmic cut and control arguments uniformly.

\section{A posteriori inverse error from a residual}
The most useful analytic bridge is a residual estimate.

\begin{theorem}[Residual-to-inverse error]\label{thm:residual-inverse}
Let $f$ be continuously differentiable and one-to-one on an interval $I$.
Suppose $x_\ast=f\compinv(y)$ and an approximation $x_N\in I$ satisfy
\[
 \AbsoluteValue{f'(x)}\ge m>0
\]
for every $x$ between $x_N$ and $x_\ast$. Then
\begin{equation}\label{eq:residual-error}
 \AbsoluteValue{x_N-x_\ast}
 \le \frac{\AbsoluteValue{f(x_N)-y}}{m}.
\end{equation}
\end{theorem}

\begin{proof}
By the mean-value theorem,
\[
 f(x_N)-f(x_\ast)=f'(\xi)(x_N-x_\ast)
\]
for some $\xi$ between $x_N$ and $x_\ast$. Take absolute values and use
$\AbsoluteValue{f'(\xi)}\ge m$.
\end{proof}

For $f(X)=X+\BigOAt{\log X}{X\to+\infty}$, one has
$f'(X)=1+\LittleO(1)$ under the relevant
differentiability hypotheses, so $m$ may be chosen close to $1$ for large
$X$. Therefore the first omitted residual block is also the natural scale
of the inverse error. This gives a rigorous reason for the residual checks
emphasized throughout the article.

\section{Newton's correction and residual squaring}
Let $x_\ast$ solve $f(x)=y$, and put
\[
 x_+=x-\frac{f(x)-y}{f'(x)}.
\]
Taylor's theorem gives
\[
 f(x_+)-y
 =\frac12f''(\xi)(x_+-x)^2
\]
for a suitable $\xi$, provided $f'$ stays nonzero. Hence a residual of
scale $\NewtonResidualScale$ is changed to scale
$\NewtonResidualScale^2$ after normalization
by a nonvanishing derivative. The formal valuation-doubling property of
Newton reversion is the algebraic shadow of this analytic estimate.

\section{Convergent, asymptotic, and divergent expansions}
Three different situations occur.
\begin{enumerate}
\item A block series may converge for sufficiently large $X$.
\item It may be a nonconvergent but valid Poincar\'e expansion, with every
      finite truncation improving the asymptotic order.
\item It may be only a formal candidate until a realization theorem is
      supplied.
\end{enumerate}

The Lambert-$\LambertW$ large-argument expansion is commonly used asymptotically,
with branch-dependent complex logarithms; its coefficients and branch
structure agree with the standard formulas in Corless et al. and the DLMF
\cite{Corless1996,DLMF413}. Convergence questions depend on how the
expansion is reorganized and on the location in the complex plane. Formal
coefficient identities should therefore not be advertised as global
convergence statements.

\section{Sectorial and parameter-uniform expansions}
For complex $X$, fix an analytic logarithm branch $\mathcal B_X$ and write
\[
 L=\LogarithmOnBranch{\mathcal B_X}{X}
\]
in a sector
\[
 \AbsoluteValue{\arg X-\theta_0}\le\theta_{\max}<\pi.
\]
A sectorial polynomial--logarithmic expansion requires remainder bounds
uniform on each closed subsector. If parameters are present, uniformity
must also be specified over a parameter set. Coefficients can change
character when a leading parameter vanishes: an invertible block may cease
to be invertible, a dominant balance may switch, or two scales may resonate.

A good theorem therefore states:
\begin{enumerate}
\item the domain and branch;
\item the parameter range;
\item the ordered scale;
\item the truncation convention;
\item the uniform remainder estimate.
\end{enumerate}

\section{Turning a formal answer into a theorem}
A standard proof strategy for an inverse expansion is:
\begin{enumerate}
\item construct a formal truncation $G_N$;
\item prove $G_N$ lies in a domain where the original function is
      one-to-one;
\item estimate the residual $R_N=f(G_N)-Y$;
\item prove a lower bound for $\AbsoluteValue{f'}$ near $G_N$;
\item apply \cref{thm:residual-inverse}.
\end{enumerate}
This method often avoids proving convergence of the full formal series.
Each finite truncation is justified directly.

\chapter{Failure modes, diagnostics, and best practices}
\label{ch:pitfalls}

Polynomial--logarithmic calculations are forgiving at leading order and
unforgiving at high order. The following failure modes account for most
incorrect formulas.

\section{Using the wrong monomial order}
For $X\to+\infty$,
\[
 X^{-1}\precdom(\log X)^{-m}\precdom1\precdom(\log X)^m\precdom X
\]
for every fixed $m>0$. Inside an $X^{-n}$ block, however, powers of
$\log X$ have their own order. Sorting only by total symbolic degree
destroys the asymptotic hierarchy.

\section{Treating \texorpdfstring{$t$}{t} and \texorpdfstring{$L$}{L} as independent under differentiation}
The formal partial derivative $\partial_t$ is not the derivative with
respect to $X$. The correct operator is
\[
 \D=t\partial_L-t^2\partial_t.
\]
Forgetting the $t\partial_L$ term loses all derivatives of logarithmic
coefficients; forgetting the $-t^2\partial_t$ term gives the wrong outer
power.

\section{Dividing before normalizing the leading block}
A quotient recurrence requires an invertible leading coefficient. First
factor the leading monomial and coefficient:
\[
 B=b_\lambda t^\lambda(1+R),\qquad \nu_tR>0.
\]
Then invert $b_\lambda$, $t^\lambda$, and $1+R$ separately. Applying a
geometric series to an expression whose ``small'' part is not actually
small is invalid.

\section{Under-expanding intermediate expressions}
To obtain an output through $t^N$, an internal logarithm, reciprocal, or
composition may need terms beyond its visually obvious order. Determine
precision by valuation dependencies, not by counting printed terms.
A final residual one block deeper is the simplest safeguard.

\section{Using Taylor composition outside its admissible regime}
The formula
\[
 F(Y+H)=\sum_{k\ge0}\frac{H^k}{k!}F^{(k)}(Y)
\]
is coefficientwise meaningful when derivatives and powers force increasing
valuation. If $H$ is comparable with or larger than $Y$, this is not a
near-identity substitution. Change the dominant scale first and use exact
generator substitution.

\section{Forgetting that composition is not commutative}
Even when $F=X+\LittleO(X)$ and $G=X+\LittleO(X)$,
\[
 F\compose G\ne G\compose F
\]
in general. The first discrepancy is governed by a differential
commutator such as $A'B-B'A$. This matters in iterated normal forms and
fractional iteration \cite{EdgarFractional}.

\section{Applying ordinary power-series reversion at a critical point}
If $F'(X_0)=0$, the inverse is generally not an ordinary integer-power
series. A zero of multiplicity $m$ leads to a Puiseux scale
$(Y-Y_0)^{1/m}$. The square-root expansion of Lambert $\LambertW$ at
$-\EulerE^{-1}$ is the canonical example. No amount of integer-power
coefficient matching repairs a wrong local scale.

\section{Suppressing branches of logarithms}
For complex $X$, fix branches $\mathcal B_1$ and $\mathcal B_2$ and write
$\LogarithmOnBranch{\mathcal B_1}{X}$ and
$\LogarithmOnBranch{\mathcal B_2}{
  \LogarithmOnBranch{\mathcal B_1}{X}}$.
Changing $\mathcal B_1$ shifts the first logarithm by an integral multiple of
$2\pi\ImaginaryUnit$ and propagates through every later block. A formula for
$\LambertW_k$ must state the branch labels and continuation domain.

\section{Writing ambiguous remainder symbols}
The estimate
\[
 \BigO(X^{-3})
\]
does not contain $X^{-3}(\log X)^{100}$ under the usual bounded-coefficient
interpretation. Prefer a scale-aware statement such as
\[
 \BigOAt{(\log X)^d/X^3}{X\to+\infty}
\]
or identify the first omitted monomial explicitly.

\section{Assuming formal differentiation preserves analytic asymptotics}
Termwise differentiation needs control of the differentiated remainder.
Oscillatory functions can have a tiny remainder and a large derivative.
Use analyticity in a sector, monotonicity plus derivative estimates, or
another theorem appropriate to the realization.

\section{Confusing equality of transseries with equality of functions}
Different functions can share the same asymptotic expansion, differing by a
term beyond every monomial in the chosen scale, such as $\EulerE^{-X}$ relative
to a pure power--log scale. Enlarging the transseries field may record that
term; staying in the smaller field deliberately forgets it.

\section{Letting coefficient domains grow unnoticed}
Starting with polynomial coefficients in $L$ does not guarantee that every
operation remains polynomial. Integration, nested-log differentiation, and
some scale changes introduce $L^{-1}$ or $\log L$. State the coefficient
ring before claiming closure.

\section{Premature floating-point evaluation}
Numerical substitution into huge logarithmic polynomials can cause
catastrophic cancellation. Keep exact coefficients, combine blocks
symbolically, and evaluate only after choosing the truncation order.
For very large arguments, evaluate in nested form rather than expanding
high-degree polynomials into monomials.

\section{No independent verification}
Deriving and simplifying an expression with the same code path is not an
independent check. Better diagnostics include:
\begin{enumerate}
\item multiply a quotient by its divisor;
\item compose an inverse in both orders;
\item compare Newton and triangular reversion;
\item verify several coefficients numerically at high precision;
\item inspect the first omitted residual block.
\end{enumerate}

\section{A practical checklist}
Before accepting a calculation, answer all of the following.

\begin{longtable}{@{}p{0.06\textwidth}p{0.86\textwidth}@{}}
\toprule
&Question\\
\midrule
1&What is the limiting variable and in which domain or sector?\\
2&What is the exact ordered monomial scale?\\
3&Which coefficient ring or iterated field is being used?\\
4&What is the leading monomial of every divisor?\\
5&Is each composition near-identity, or has the scale first been changed?\\
6&Which logarithmic branches are fixed?\\
7&What target filtration and guard precision are required?\\
8&Can discarded terms feed a retained block?\\
9&Does the derivative operator respect $t=X^{-1}$ and $L=\log X$?\\
10&Does the inverse map have nonzero leading derivative?\\
11&Is the result formal, asymptotic, or convergent?\\
12&What analytic hypotheses justify composition or differentiation?\\
13&Has the result been substituted back into the defining equation?\\
14&Does the first nonzero residual have the predicted order?\\
15&Have coefficients been kept exact until final evaluation?\\
\bottomrule
\end{longtable}

\begin{summarybox}
The safest philosophy is filtration first, algebra second, simplification
third. Choose the scale, prove that every operation is coefficientwise
finite, truncate aggressively, and certify the answer by a residual.
\end{summarybox}


\appendix

\chapter{Bell polynomials and Fa\`a di Bruno bookkeeping}
\label{app:bell}

Bell polynomials organize the derivative combinatorics behind composition,
Taylor substitution, and Lagrange inversion.

\section{Partial exponential Bell polynomials}
For integers $n\ge k\ge0$, define
\begin{equation}\label{eq:bell-definition}
 B_{n,k}(x_1,x_2,\ldots)
 =
 \sum
 \frac{n!}{j_1!j_2!\cdots}
 \prod_{m\ge1}\left(\frac{x_m}{m!}\right)^{j_m},
\end{equation}
where the sum is over all nonnegative integer sequences satisfying
\[
 \sum_{m\ge1}j_m=k,\qquad
 \sum_{m\ge1}m j_m=n.
\]
The complete Bell polynomial is
\[
 B_n(x_1,\ldots,x_n)=\sum_{k=0}^nB_{n,k}(x_1,\ldots,x_{n-k+1}).
\]

The first partial polynomials include
\begin{align*}
 B_{1,1}&=x_1,\\
 B_{2,1}&=x_2,&B_{2,2}&=x_1^2,\\
 B_{3,1}&=x_3,&B_{3,2}&=3x_1x_2,&B_{3,3}&=x_1^3,\\
 B_{4,2}&=4x_1x_3+3x_2^2,&
 B_{4,3}&=6x_1^2x_2,&B_{4,4}&=x_1^4.
\end{align*}

\section{Fa\`a di Bruno's formula}
For sufficiently differentiable functions,
\begin{equation}\label{eq:faa-di-bruno}
 \D^n(f\compose g)
 =
 \sum_{k=0}^{n}
 f^{(k)}(g)\,
 B_{n,k}\bigl(g',g'',\ldots,g^{(n-k+1)}\bigr).
\end{equation}
In formal differential algebras this is an algebraic identity. It explains
why the $n$th composition coefficient depends polynomially on derivatives
of the source coefficients and on partitions of $n$.

\begin{proof}
Expand
\[
 f(g+h)=\sum_{k\ge0}\frac{f^{(k)}(g)}{k!}
       (g(X+h)-g(X))^k.
\]
Now
\[
 g(X+h)-g(X)
 =\sum_{m\ge1}\frac{g^{(m)}(X)}{m!}h^m.
\]
The coefficient of $h^n/n!$ in the $k$th power is exactly
$B_{n,k}(g',g'',\ldots)$. Comparing coefficients proves
\eqref{eq:faa-di-bruno}.
\end{proof}

\section{Bell-polynomial form of implicit coefficients}
Suppose $G=Y+H$ and $F=X+A$. The inverse equation is
\[
 H=-A(Y+H).
\]
Taylor expansion gives
\[
 H=-\sum_{k\ge0}\frac{H^k}{k!}A^{(k)}(Y).
\]
Equating outer blocks recursively produces polynomials in
$A,A',A'',\ldots$. The first universal terms are
\begin{align}
 H={}&-A+AA'
 -A(A')^2-\frac12A^2A''\notag\\
 &+A(A')^3+\frac32A^2A'A''
 +\frac16A^3A'''+\cdots.
 \label{eq:universal-inverse-derivatives}
\end{align}
These are the expansion of
\[
 \sum_{n\ge1}\frac{(-1)^n}{n!}\D^{n-1}(A^n).
\]
Bell polynomials provide an efficient way to expand the repeated
derivatives without repeatedly applying the product rule from scratch.

\section{Coefficient extraction with an auxiliary variable}
A particularly clean implementation introduces an auxiliary ordinary
variable $z$ and uses
\[
 \exp\!\left(\sum_{m\ge1}x_m\frac{z^m}{m!}\right)
 =\sum_{n\ge0}B_n(x_1,\ldots,x_n)\frac{z^n}{n!}.
\]
All Bell polynomials up to degree $N$ can then be generated by a single
truncated exponential. This technique is useful when composition
coefficients are needed symbolically in terms of unspecified input blocks.

\chapter{Changing between infinity and zero}
\label{app:zero-infinity}

Many power--log expansions are written at $x\to0^+$ rather than
$X\to+\infty$. The conversion is elementary but sign errors are common.

\section{The basic dictionary}
Set
\[
 X=x^{-1},\qquad L=\log X=-\log x.
\]
Then
\begin{equation}\label{eq:zero-infinity-dictionary}
 X^{-n}P(\log X)
 =x^nP(-\log x).
\end{equation}
Thus
\[
 \sum_{n\ge n_0}X^{-n}p_n(\log X)
 \quad\longleftrightarrow\quad
 \sum_{n\ge n_0}x^np_n(-\log x).
\]
If $\ell=\log L$, then
\[
 \ell=\log(-\log x)
\]
on the positive real axis.

\section{Derivative conversion}
Since $X=x^{-1}$,
\[
 \frac{\Differential}{\Differential x}=-X^2\frac{\Differential}{\Differential X}.
\]
In the zero-side variables $x$ and $\lambda=-\log x$,
\[
 \frac{\Differential}{\Differential x}
 =\partial_x-\frac1x\partial_\lambda
\]
when $x$ and $\lambda$ are temporarily treated as formal generators.
Applying this to $x^nP(\lambda)$ gives
\[
 \frac{\Differential}{\Differential x}\bigl(x^nP(\lambda)\bigr)
 =x^{n-1}\bigl(nP(\lambda)-P'(\lambda)\bigr),
\]
the exact counterpart of \eqref{eq:implementation-D-block}.

\section{The lower Lambert branch}
For $x\to0^+$, the solution $\LambertW_{-1}(-x)$ tends to $-\infty$. Put
\[
 S=\log(1/x),\qquad M=\log S.
\]
Writing $\LambertW_{-1}(-x)=-U$ transforms
\[
 \LambertW_{-1}(-x)\EulerE^{\LambertW_{-1}(-x)}=-x
\]
into
\[
 U-\log U=S.
\]
Reversion on the appropriate sign branch gives
\[
 \LambertW_{-1}(-x)
 =-S-M-\frac{M}{S}
 +\frac{M^2/2-M}{S^2}
 +\BigOAt{M^3/S^3}{x\to0^+}.
\]
The same Lambert polynomials occur, but signs and the logarithmic branch are
changed by the transformed defining equation.

\section{Dulac-style power--log series}
Near $x=0^+$, a finite-type power--log series often has the form
\[
 \sum_{\alpha\in A}x^\alpha p_\alpha(\log x),
\]
where $A$ is a well-ordered additive subset of positive reals. Such series
appear in local dynamical normal forms and Dulac maps. The power--log
literature develops composition, normal forms, and embeddings in this local
language \cite{MRRZ2016}. Formula \eqref{eq:zero-infinity-dictionary}
shows the direct relation with the large-variable conventions of this
article, while the admissible supports may be more general than consecutive
integer powers.

\chapter{Quick-reference formula sheet}
\label{app:quick-reference}

This appendix collects the main formulas without proofs.

\section*{Variables and scale}
\[
 t=X^{-1},\qquad L=\log X,\qquad
 \PLpoly=\K[L]((t)),\qquad
 \PLfield=\K((L^{-1}))((t)).
\]
For $X\to+\infty$,
\[
 X^{-1}\precdom L^{-m}\precdom1\precdom L^m\precdom X
 \qquad(m>0).
\]

\section*{Multiplication}
If
\[
 A=\sum A_nt^n,\qquad B=\sum B_nt^n,
\]
then
\[
 AB=\sum_rC_rt^r,\qquad
 C_r=\sum_{i+j=r}A_iB_j.
\]

\section*{Reciprocal and quotient}
For $B=\sum_{n\ge0}b_nt^n$ with invertible $b_0$,
\[
 c_0=b_0^{-1},\qquad
 c_n=-b_0^{-1}\sum_{j=1}^nb_jc_{n-j},
 \qquad B^{-1}=\sum c_nt^n.
\]
For $Q=A/B$,
\[
 q_n=b_0^{-1}\left(a_n-\sum_{j=1}^nb_jq_{n-j}\right).
\]

\section*{Powers, logarithm, and exponential}
For $\nu_tU>0$,
\begin{align*}
 (1+U)^\alpha&=\sum_{k\ge0}\binom{\alpha}{k}U^k,\\
 \log(1+U)&=\sum_{k\ge1}\frac{(-1)^{k+1}}{k}U^k,\\
 \exp U&=\sum_{k\ge0}\frac{U^k}{k!}.
\end{align*}

\section*{Differentiation}
\[
 \D=t\partial_L-t^2\partial_t,
\]
\[
 \D(P(L)t^n)=\bigl(P'-nP\bigr)t^{n+1},
\]
\[
 \D^m(P(L)t^n)
 =t^{n+m}\prod_{j=0}^{m-1}(\partial_L-n-j)P(L).
\]

\section*{Near-identity composition}
For $G(Y)=Y(1+U)$,
\[
 t_X\mapsto t_Y(1+U)^{-1},
 \qquad
 L_X\mapsto L_Y+\log(1+U).
\]
Equivalently, for coefficientwise summable Taylor families,
\[
 F(Y+H)=\sum_{k\ge0}\frac{H^k}{k!}F^{(k)}(Y).
\]

\section*{Near-identity reversion}
For $F(X)=X+A(X)$,
\begin{equation}\label{eq:quick-LB}
 F\compinv(Y)
 =Y+\sum_{n\ge1}\frac{(-1)^n}{n!}
   \D^{n-1}\!\left(A(Y)^n\right).
\end{equation}
The first terms are
\[
 Y-A+AA'
 -A(A')^2-\frac12A^2A''
 +A(A')^3+\frac32A^2A'A''
 +\frac16A^3A'''+\cdots.
\]
Newton iteration is
\[
 G_{\mathrm{new}}
 =G-\frac{F(G)-Y}{F'(G)}.
\]

\section*{Generic affine-log inverse}
If
\[
 F(X)=X+a\log X+b
 +\frac{c\log X+d}{X}
 +\BigOAt{(\log X)^2/X^2}{X\to+\infty},
\]
then
\[
 F\compinv(Y)
 =Y-aL-b
 +\frac{(a^2-c)L+ab-d}{Y}
 +\BigOAt{L^2/Y^2}{Y\to+\infty}.
\]

\section*{Lambert and Wright omega}
The Wright omega function is the inverse of $X+\log X$:
\[
 \WrightOmega(Y)+\LogarithmOnBranch{\mathcal B_\WrightOmega}{\WrightOmega(Y)}=Y,\qquad
 \LambertW_k(z)=\WrightOmega\!\left(\PrincipalLogarithm z+2\pi\ImaginaryUnit k\right)
\]
with compatible branches. Its expansion is
\[
 \WrightOmega(Y)=Y-L+\sum_{r\ge1}\frac{\LambertPolynomial{r}(L)}{Y^r},
 \qquad L=\log Y,
\]
where
\begin{align*}
 \LambertPolynomial{r}(z)
 &=\frac{(-1)^{r+1}}{(r+1)!}
   \prod_{j=0}^{r-1}(\partial_z-j)z^{r+1}\\
 &=\sum_{m=1}^r
   \frac{(-1)^{r-m}}{m!}
   \stirlingone{r}{r-m+1}z^m.
\end{align*}
The first blocks are
\[
 \LambertPolynomial{1}=z,\qquad
 \LambertPolynomial{2}=\frac{z^2}{2}-z,\qquad
 \LambertPolynomial{3}=\frac{z^3}{3}-\frac32z^2+z.
\]
Thus, with $L_1=\log z$ and $L_2=\log L_1$,
\[
 \LambertW(z)=L_1-L_2+\frac{L_2}{L_1}
 +\frac{L_2^2/2-L_2}{L_1^2}
 +\frac{L_2^3/3-\frac32L_2^2+L_2}{L_1^3}
 +\cdots.
\]

\section*{Residual certification}
If $x_N$ approximates $f\compinv(y)$ and
$\AbsoluteValue{f'}\ge m>0$ between $x_N$ and the exact inverse, then
\[
 \AbsoluteValue{x_N-f\compinv(y)}
 \le\frac{\AbsoluteValue{f(x_N)-y}}{m}.
\]

\chapter{Notation and terminology}
\label{app:notation}

\begin{longtable}{@{}p{0.20\textwidth}p{0.72\textwidth}@{}}
\toprule
Symbol or term&Meaning\\
\midrule
$X$&Large independent variable, normally tending to $+\infty$.\\
$t=X^{-1}$&Outer small generator.\\
$L=\log X$&First logarithmic generator.\\
$\log_2X$&Iterated logarithm $\log\log X$ when defined.\\
$\K$&Characteristic-zero coefficient field, commonly $\RealNumbers$ or $\ComplexNumbers$.\\
$\PLpoly$&Polynomial-coefficient Laurent algebra $\K[L]((t))$.\\
$\PLfield$&Iterated Laurent completion, typically
$\K((L^{-1}))((t))$.\\
$\nu_tT$&Least outer exponent with nonzero coefficient.\\
$\CoefficientSupportOf{T}$&Set of monomials or outer exponents occurring in $T$.\\
$\precdom$&Strict asymptotic dominance: $f/g\to0$.\\
$\compose$&Composition.\\
$F\compinv$&Compositional inverse of $F$.\\
$\D$&Derivative with respect to the large variable.\\
block&All monomials sharing one outer power $t^n$.\\
filtration&Nested additive tail subgroups used to define precision and
coefficientwise convergence; after restriction to a power-series tail ring,
these subgroups are ideals.\\
near identity&A map $F(X)=X+\LittleO(X)$ in the selected scale.\\
series reversal&Compositional inversion, not reciprocal.\\
summable family&A formal family for which every monomial receives only
finitely many contributions.\\
grid support&Support generated by finitely many small ratios, with
well-founded exponent restrictions.\\
residual&The defining equation evaluated at a truncated approximation.\\
\bottomrule
\end{longtable}

\backmatter

\chapter*{Concluding perspective}
\addcontentsline{toc}{chapter}{Concluding perspective}
Polynomial--logarithmic transseries occupy a productive middle ground.
They retain the ordered-support discipline of general transseries while
remaining close enough to ordinary Laurent series for explicit
coefficient algorithms. Their apparent complication is concentrated in
two facts: logarithms form subordinate scales inside power blocks, and
composition moves both the inverse-power generator and the logarithmic
generator.

Once those facts are encoded, the main operations become systematic.
Multiplication is convolution. Division is triangular. Near-identity
composition is either a two-generator substitution or a filtered Taylor
series. Reversal is a triangular group operation, accelerated by Newton
iteration and summarized in closed form by Lagrange--B\"urmann. Lambert
$\LambertW$ is the canonical example because its logarithmic defining equation is
exactly the inverse of $X+\log X$, but the same machinery governs a much
wider range of inverse asymptotics.

The most durable computational lesson is to organize every calculation by
valuation and to verify every claimed truncation by substitution. The most
important conceptual lesson is to separate formal algebra from analytic
realization. Formal transseries determine the only possible coefficients;
residual estimates, monotonicity, sectorial bounds, or realization theorems
then decide when those coefficients describe an actual function.

\begin{thebibliography}{99}

\bibitem{Corless1996}
R.~M. Corless, G.~H. Gonnet, D.~E.~G. Hare, D.~J. Jeffrey, and
D.~E. Knuth,
\newblock On the Lambert \(\LambertW\) function,
\newblock \emph{Advances in Computational Mathematics} \textbf{5} (1996),
329--359.
\newblock \url{https://doi.org/10.1007/BF02124750}.

\bibitem{DLMF413}
NIST Digital Library of Mathematical Functions,
\newblock \S4.13, Lambert \(\LambertW\)-function,
\newblock Version 1.2.7, release date June 15, 2026,
accessed August 31, 2026.
\newblock \url{https://dlmf.nist.gov/4.13}.

\bibitem{EdgarBeginners}
G.~A. Edgar,
\newblock Transseries for beginners,
\newblock \emph{Real Analysis Exchange} \textbf{35} (2009/2010), no.~2,
253--310.
\newblock \url{https://arxiv.org/abs/0801.4877}.

\bibitem{EdgarComposition}
G.~A. Edgar,
\newblock Transseries: composition, recursion, and convergence,
\newblock arXiv:0909.1259, 2009.
\newblock \url{https://arxiv.org/abs/0909.1259}.

\bibitem{EdgarFractional}
G.~A. Edgar,
\newblock Fractional iteration of series and transseries,
\newblock \emph{Transactions of the American Mathematical Society}
\textbf{365} (2013), 5805--5832.
\newblock \url{https://arxiv.org/abs/1002.2378}.


\bibitem{KaluginJeffrey2012}
G.~A. Kalugin and D.~J. Jeffrey,
\newblock Convergence in $\ComplexNumbers$ of series for the Lambert $\LambertW$ function,
\newblock arXiv:1208.0754, 2012.
\newblock \url{https://arxiv.org/abs/1208.0754}.

\bibitem{MRRZ2016}
P.~Marde\v{s}i\'c, M.~Resman, J.-P. Rolin, and V.~\v{Z}upanovi\'c,
\newblock Normal forms and embeddings for power-log transseries,
\newblock \emph{Advances in Mathematics} \textbf{303} (2016), 888--953.
\newblock \url{https://doi.org/10.1016/j.aim.2016.08.030}.

\bibitem{SalvyShackell1999}
B.~Salvy and J.~Shackell,
\newblock Symbolic asymptotics: multiseries of inverse functions,
\newblock \emph{Journal of Symbolic Computation} \textbf{27} (1999),
no.~6, 543--563.
\newblock \url{https://doi.org/10.1006/jsco.1999.0281}.

\bibitem{Shackell2004}
J.~Shackell,
\newblock \emph{Symbolic Asymptotics},
\newblock Springer, Berlin, 2004.

\bibitem{vanDenDries2001}
L.~van den Dries, A.~Macintyre, and D.~Marker,
\newblock Logarithmic-exponential series,
\newblock \emph{Annals of Pure and Applied Logic} \textbf{111} (2001),
61--113.
\newblock \url{https://doi.org/10.1016/S0168-0072(01)00035-5}.

\bibitem{vdHoeven2006}
J.~van der Hoeven,
\newblock \emph{Transseries and Real Differential Algebra},
\newblock Lecture Notes in Mathematics, vol.~1888,
Springer, Berlin, 2006.
\newblock \url{https://doi.org/10.1007/3-540-35590-1}.

\end{thebibliography}

\end{document}
